\chapter{个性化排序与响应预估}
\label{chap:rs-ranking}

召回已经把数百万件物品缩小为数百个候选，但登山包、水壶和相机之间仍需作出选择。用户最可能点击的物品，未必带来最高的消费价值；购买概率最高的商品，也可能在没有推荐时就会被购买。评分器必须先回答“分数代表什么”，才能讨论用什么网络计算它。

本章以给定候选后的比较为中心。响应与价值预估描述可能发生的结果，偏好排序学习对象之间的次序，增量排序寻找推荐能改变的结果。确定目标后，字段交互、历史行为和语义内容提供不同的计算依据；这些依据既能由分离模块融合，也能进入联合骨干。日志怎样形成、分数怎样使用与维护，则决定了学到的关系能否支持实际比较。最终展示还涉及对象之间的关系，留到下一章展开。

% Outline: 4.1
\section{预估与排序任务}
\label{sec:rs-ranking-task}

设一次请求的可用信息为$\bm{x}_t$，其中只含时点$t$之前可取得的用户资料、历史$H_{u,t}$及情境；候选集合为$C_t$，候选$i$的属性可以并入评分输入。评分函数$s_\theta(\bm{x}_t,i)$输出一个可比较的数。如果它通过sigmoid形成$p_\theta=\sigma(s_\theta)$，并按点击标签训练，才有了点击概率模型的形式；是否校准还需数据检验。任意单调变换都保持排序次序，却会改变概率解释。

响应必须连同条件和窗口定义。某商品“24小时内购买”的概率，与“曝光后经点击并在7天内购买”的概率不是同一目标。展示位置、页面、价格和库存也会影响响应。排序时尚未确定的位置不能随意填入一个值，再把输出解释成脱离位置的偏好。若任务面向家庭或群组，评分对象还要从个人改为成员集合；第~\ref{sec:rs-agree}、\ref{sec:rs-groupim}~节给出的群组表示可以复用，但成员满意度与群体决策规则仍需分别定义。

预测、排序与决策因此有三种不同的错误。把所有点击概率高估一倍可能不改变次序，却使按期望收益分配资源的决策失准；把前两名交换可能只略微增加平均预测误差，却影响首位展示；准确预测用户喜欢某件物品，也不能证明重复展示十次值得。下文先把这些目标分开，再建立共同评分器。

% Outline: 4.2
\section{预估与排序目标}
\label{sec:rs-ranking-objectives}

选择损失函数之前，应先写出要估计的量及其标签来源。相同的矩阵分解或深层网络，可以拟合显式评分、学习点击排序，也可以在有处理与对照证据时学习推荐增量。架构名称不能替代目标定义。

% Outline: 4.2.1
\subsection{响应与价值预估}
\label{sec:rs-response-value}

令$Y$为某一明确窗口内的二值响应，$V$为响应发生后的价值。把未发生响应时的价值定义为零，则
\begin{equation}
 \mathbb E[YV\mid\bm x,i]
 =\mathbb P(Y=1\mid\bm x,i)\,
   \mathbb E[V\mid Y=1,\bm x,i].
 \label{eq:rs-response-value}
\end{equation}
这个等式来自条件期望，不要求响应与价值独立。教学例中，背包的购买概率为0.04，购买后平均毛利为80元，期望毛利为3.2元；水壶的相应值为0.12和20元，期望毛利为2.4元。购买概率排序选择水壶，期望毛利排序选择背包，二者回答的问题不同。这里的毛利还没有扣除推荐成本，也没有识别增量。

二值响应常用交叉熵$\ell_{\mathrm{BCE}}(y,p)=-y\log p-(1-y)\log(1-p)$；数值价值可用平方损失拟合条件均值，也可按分布和应用选择其他目标。观看时长的零值、截断和长尾应在标签中定义，满意度调查则只覆盖被邀请且实际回答的人，不能把缺失回答填成“不满意”。

% Outline: 4.2.1.1
\subsubsection{多行为联合预估与共享基线}
\label{sec:rs-shared-bottom}

一个视频候选可能同时需要点击、观看和满意度预测。共享底层（shared-bottom）用同一表示网络$f_\theta$处理输入，再由任务头$h_k$输出$\widehat y_k=h_k(f_\theta(\bm x,i))$。它让多个任务共同更新底层参数，而单任务基线为每个目标独立训练一套网络。YouTube多任务排序研究同时使用参与和满意度反馈，恰好说明“容易记录的行为”与“想要了解的体验”并不等价\scite{rs-r16}

设样本$n$在任务$k$的标签可用标记为$m_{n,k}\in\{0,1\}$。一种明确处理缺失的联合目标是
\begin{equation}
 \mathcal L=\sum_{k=1}^{K}\lambda_k
 \frac{\sum_n m_{n,k}\ell_k(y_{n,k},\widehat y_{n,k})}
      {\sum_n m_{n,k}},
 \label{eq:rs-multitask-loss}
\end{equation}
其中只包含当前批次有标签的任务，$\lambda_k$调节训练贡献。按任务分别归一化，可以避免标签多的任务仅凭数量占据全部梯度，但仍需关注损失尺度与采样分布。

例如，一条记录显示点击为1、观看时长为35秒，却没有满意度问卷。若点击预测0.7、时长预测30秒，点击损失为$-\log0.7\approx0.357$；以$\tfrac12(y-\widehat y)^2$为时长损失，则该项为12.5。满意度项被掩码排除。未点击样本的时长是否为零，取决于预测的是“曝光后总观看时长”还是“点击后观看时长”，不能只靠掩码习惯决定。

共享是否有益，应比较相同数据、预算下的各单任务误差。若点击改善而满意度退化，可能存在负迁移，也可能只是标签覆盖或权重变化；应分组检查。式~\eqref{eq:rs-multitask-loss}的权重服务训练，不自动等于最终展示效用的权重。

% Outline: 4.2.1.2
\subsubsection{MMoE：任务门控响应预估}
\label{sec:rs-mmoe}

共享底层让所有任务固定读取同一个表示，可能迫使点击与满意度共享不合适的信息；单任务又各自训练、无法共用输入编码。\term{多门控混合专家}（multi-gate mixture-of-experts, MMoE）保留一组共享专家，但让每个任务决定如何混合专家输出\scite{rs-r37}

\textbf{模型结构。}第一步，同一候选输入$\bm z$并行送入$J$个专家，各自编码为$\bm e_j=f_j(\bm z)$，专家参数跨任务共享；第二步，任务$k$的门控以$\bm z$为输入softmax出混合权重；第三步，按权重混合专家得到任务表示$\bm h_k$，再由任务头输出预测：
\begin{equation}
 \bm g_k=\operatorname{softmax}(\bm W_k\bm z),\qquad
 \bm h_k=\sum_{j=1}^{J}g_{k,j}\bm e_j,\qquad
 \widehat y_k=h_k(\bm h_k).
 \label{eq:rs-mmoe}
\end{equation}
专家由多个任务共享，门控及预测头按任务区分；各响应标签仍使用式~\eqref{eq:rs-multitask-loss}，只纳入本批有标签的任务。通用专家网络的容量与路由问题不改变这些标签的含义。

教学例取三个二维专家输出$(1,0),(0,1),(1,1)$。点击门控为$(0.6,0.1,0.3)$，得到$(0.9,0.4)$；满意度门控为$(0.1,0.7,0.2)$，得到$(0.3,0.9)$。若两个教学预测头分别采用$\sigma(h_1-h_2)$与$\sigma(h_2-h_1)$，输出约为0.622和0.646。同样的专家能向两个任务提供不同组合，而任务头仍可进一步变换。

训练时所有专家都接受各任务回传的梯度，门控按任务区分路由；服务时软混合通常仍计算全部专家，不能把较小门控权重直接算成省掉的运算。门控相近可以帮助诊断共享情况，却不能证明点击造成满意度或两个行为具有同一因果机制。是否减轻负迁移，仍应逐任务检验并计入专家计算。

% Outline: 4.2.1.3
\subsubsection{PLE：共享与任务专属响应预估}
\label{sec:rs-ple}

MMoE的专家都向各任务开放；若某些信息只对观看有用，另一任务仍可能通过共享专家把无关影响带进来。\term{渐进分层提取}（progressive layered extraction, PLE）为每个任务保留专属专家，另设共享专家，再逐层交换和提取信息\scite{rs-r38}

\textbf{模型结构。}第一步，每层同时维护一组共享专家与每个任务自己的专属专家；第二步，任务$k$的门控只能混合本任务专家与共享专家，共享门控则可混合各任务及共享专家；第三步，得到的任务表示进入下一层的本任务专家，共享表示进入下一层的共享专家，逐层传递；第四步，两层以后任务头读取各自最终表示。较早的共享通路可以传递跨任务知识，较后的任务通路逐渐形成分工；“专属”不意味着这些参数在整个计算图中完全不受其他任务间接影响。

用观看和满意度两个任务作标量教学例。第一层三个专家输出分别为观看2、满意度$-1$、共享1；任务门控在“自身、共享”上分别取$(0.75,0.25)$与$(0.5,0.5)$，产生1.75与0；共享门控取$(0.2,0.3,0.5)$，产生0.6。若第二层的三个专家暂取恒等映射，观看门控$(0.6,0.4)$产生1.29，满意度门控$(0.4,0.6)$产生0.36。再接各自预测头，即可按已有标签联合训练。真实专家当然可以是多层非线性网络。

训练按任务掩码损失联合更新专属专家、共享专家、门控与任务头；服务时各任务门控只在允许的专家集合内混合，专属专家不向其他任务开放。这个计算展示了门控可见范围，没有证明任务都得到改善。检查PLE时，应同时报告观看和满意度相对单任务、共享底层及MMoE的变化，并固定样本覆盖和预算。多个准确的响应头仍需由第~\ref{chap:rs-constraints}~章的目标与约束决定如何使用。

% Outline: 4.2.1.4
\subsubsection{NMTR：级联行为的联合预估}
\label{sec:rs-nmtr}

MMoE、PLE把各行为当作并行任务，没有利用浏览、加购、购买在业务上规定的先后顺序。\term{神经多任务推荐}（neural multi-task recommendation, NMTR）共享用户、物品嵌入，为各行为保留交互函数，再把前一行为的预测加到后一行为的logit中\scite{rs-r97}

\textbf{模型结构。}第一步，用户与物品嵌入在三种行为间共享；第二步，每个行为有自己的交互函数$f_k(\bm h_u,\bm v_i)$；第三步，按业务顺序把前一行为的预测$\widehat y_{k-1}$接入后一行为logit。以三种行为为例，
\begin{align}
 \widehat y_1&=\sigma(f_1(\bm h_u,\bm v_i)+b_1),\\
 \widehat y_2&=\sigma(f_2(\bm h_u,\bm v_i)+\widehat y_1+b_2),\\
 \widehat y_3&=\sigma(f_3(\bm h_u,\bm v_i)+\widehat y_2+b_3).
 \label{eq:rs-nmtr}
\end{align}
各行为的已观测正例及采样未观测对象形成各自交叉熵，联合更新共享嵌入和行为模块；服务时按顺序计算级联输出，以指定目标行为的分数比较候选。

若教学浏览预测为0.6，加购交互logit为$-1$且偏置为零，则加购预测为$\sigma(-0.4)\approx0.401$。若该logit改为1，加购预测会达到$\sigma(1.6)\approx0.832$，高于浏览预测。因此NMTR的级联是预测连接，不是条件概率相乘，也不施加行为概率单调约束。未观测对象仍是带假设的训练负例。

行为顺序应符合当前应用的定义。不能为了凑齐级联而把浏览复制成购买标签；跨设备或其他入口也可能使记录顺序不完整。第~\ref{sec:rs-mbgcn}~节的多行为图利用行为作为输入证据，本节利用多种行为提供监督；广告中的点击后转化概率链路则见第~\ref{chap:rs-advertising}~章。

% Outline: 4.2.2
\subsection{偏好与排序学习}
\label{sec:rs-ranking-learning}

如果目标是把更合适的对象放在前面，不必先得到精确的响应概率。逐对象目标分别拟合每件物品的标签；成对目标直接比较$i$与$j$的分差；列表目标在同一候选列表内分配概率或惩罚次序。例如，以人工相关性$r_i$构造目标分布$q_i$，可最小化$-\sum_{i\in C}q_i\log\operatorname{softmax}(\bm s)_i$。这个归一化只在候选$C$内成立，不等于单件物品被点击的概率。

比较单位改变了训练强调的错误。逐对象损失关心数值偏离，成对损失关心相对关系，列表损失可以强调头部竞争。三者都可以配合不同评分架构，也都要说明候选怎样取得、标签如何观察。列表排序损失仍不直接描述一组物品之间的互补、重复或展示成本。

% Outline: 4.2.2.1
\subsubsection{BPR：隐式反馈成对排序}
\label{sec:rs-bpr}

逐对象拟合把未观测物品一律当负例，可隐式反馈里只有“已交互优于采样到的未观测”这一相对关系更可靠。\term{贝叶斯个性化排序}（Bayesian personalized ranking, BPR）把这一假设写成三元组$(u,i,j)$：$i$已交互，$j$为采样未观测\scite{rs-h04}

\textbf{模型结构。}第一步，按$(u,i,j)$三元组采样，$i$已交互、$j$为采样未观测；第二步，计算分差$\Delta_{u,i,j}=s(u,i)-s(u,j)$，用$\sigma(\Delta_{u,i,j})$建模$i$优于$j$的偏好关系；第三步，结合零均值高斯参数先验，得到负对数后验目标
\begin{equation}
 \mathcal L_{\mathrm{BPR}}
 =-\sum_{(u,i,j)\in D}\log\sigma(\Delta_{u,i,j})
  +\frac{\lambda}{2}\lVert\bm\theta\rVert_2^2.
 \label{eq:rs-bpr}
\end{equation}
对三元组采样作随机梯度更新即可训练。分差越小，正负次序越不确定，$-\log\sigma(\Delta)$越大。

以第~\ref{sec:rs-mf}~节的MF评分为例，$\Delta=\bm h_u^\top(\bm v_i-\bm v_j)$。忽略正则的单步教学更新取$\bm h_u=(0.6,0.2)$、$\bm v_i=(0.4,0.5)$、$\bm v_j=(0.2,0.1)$，得$\Delta=0.2$、损失约0.598。记$g=\sigma(-\Delta)\approx0.4502$，学习率0.1，则同时使用旧向量更新：
\begin{align}
 \bm h'_u&=\bm h_u+0.1g(\bm v_i-\bm v_j)
           \approx(0.6090,0.2180),\\
 \bm v'_i&=\bm v_i+0.1g\bm h_u\approx(0.4270,0.5090),\\
 \bm v'_j&=\bm v_j-0.1g\bm h_u\approx(0.1730,0.0910).
\end{align}
更新后的分差约为0.2458，正物品相对优势增大。加入正则时还应减去对应的参数收缩项。

加权隐式MF拟合二值偏好及置信度，BPR-MF拟合成对关系；相同因子结构不意味着相同目标。BPR也能训练LightGCN等评分器。未观测$j$可能只是没有被展示，故“三元组为真”是学习假设；均匀负采样与热门负采样又赋予不同的比较权重。输出按次序使用，并以相同候选上的排名和截断指标评价，不要求$s$具有概率含义。

% Outline: 4.2.3
\subsection{因果增量与排序学习}
\label{sec:rs-uplift-ranking}

高响应并不等于高增量。设$T\in\{0,1\}$表示是否实施已经定义好的推荐处理，$Y(1),Y(0)$为相应潜在结果。条件平均处理效应
\begin{equation}
 \tau(\bm x,i)=\mathbb E[Y(1)-Y(0)\mid\bm x,i]
 \label{eq:rs-uplift}
\end{equation}
描述同类机会中推荐造成的平均变化。处理可以是“在约定位置展示”，也可以是某个完整推荐方案；改变处理定义就改变了$\tau$。每次机会只能观察一个潜在结果，因此需要对照证据和识别条件，不能把普通响应网络的输出改名为增量。

% Outline: 4.2.3.1
\subsubsection{DLCE：推荐增量的排序学习}
\label{sec:rs-dlce}

按展示后购买率排序可能选中本来就会买的物品，普通响应网络的输出不能改名为增量。教学例中，商品A展示与不展示时的购买率分别为0.9、0.85，商品B为0.6、0.3。按展示后的购买率，A领先；按增量，B的0.3高于A的0.05。\term{DLCE}直接学习这种增量的排序，用推荐处理、购买结果及处理倾向构造因果排名目标\scite{rs-r121}

\textbf{模型结构。}第一步，令$e(\bm x,i)=\mathbb P(T=1\mid\bm x,i)$为处理倾向，在一致性、给定特征后无未观测混杂及$0<e<1$等条件下，用逆倾向加权构造观测贡献
\begin{equation}
 \widehat\tau=\frac{TY}{e}-\frac{(1-T)Y}{1-e}
 \label{eq:rs-uplift-ips}
\end{equation}
其条件期望为$\tau$，它可以取较大的正负值，不是单次机会的真实个体效应。第二步，因果平均名次把物品名次乘以其效应后平均，期望让正增量对象取得较小名次、负增量对象取得较大名次；将名次展开为候选间的比较，即可形成可训练的排序目标。第三步，DLCE分别保留处理正贡献与对照负贡献，不能直接把带符号的$\widehat\tau$随意乘到普通正权重排序上界：负权重会改变不等式方向。设$\Delta=s(u,i)-s(u,j)$，将对照项的负名次改写为反向比较及与参数无关的常数，得到两种方向的logistic代理：
\begin{equation}
 \ell_{\mathrm{DLCE}}=
 \frac{TY}{\max(e,\chi_1)}\log(1+\exp(-\omega\Delta))
 +\frac{(1-T)Y}{\max(1-e,\chi_0)}
       \log(1+\exp(\omega\Delta)).
 \label{eq:rs-dlce-loss}
\end{equation}
这里$\omega>0$控制平滑程度，$\chi_1,\chi_0$为倾向截断阈值。若需把自然对数形式严格写成0--1比较的上界，可整体除以$\log2$；损失的两种方向不变。处理购买推动$i$上升，对照购买推动$i$下降。原方法从购买记录采样$(u,i)$，再采样比较物品$j$，以MF和正则化随机梯度训练。输出用于增量次序比较，不要求等于校准的$\tau$。

例如处理概率为0.2的一次购买记录，处理项贡献为5；对照概率为0.8的一次购买记录，对照项贡献为$-1.25$。极小概率会放大方差，截断倾向或权重又会引入偏差。若未展示时根本无法追踪购买，就不能把未记录当成$Y(0)=0$。位置纠偏试图恢复响应学习依据，DLCE进一步改变要优化的目标；整体增量实验由第~\ref{chap:rs-evaluation}~章建立。

% Outline: 4.3
\section{候选评分方案}
\label{sec:rs-scoring}

响应、偏好和增量目标都需要一个分数，但输入依据与组织方式可以不同。本节先固定候选，比较字段、历史和语义的处理，再研究它们如何构成完整评分器。观测学习和应用维护是这些方案共同面对的条件，不是另外一种与神经网络并列的模型分类。

% Outline: 4.3.1
\subsection{评分依据与处理方法}
\label{sec:rs-scoring-evidence}

同一背包在工作日通勤与周末登山中可能得到不同响应。字段交互刻画用户、对象及情境的组合；历史评分寻找与候选有关的行为证据；语义评分利用内容及需求的含义。三者可以同时出现在一个模型里。下面比较方法时，默认固定标签窗口、候选范围与请求时可用信息，避免把多读了一类输入误当成计算结构的优势。

% Outline: 4.3.1.1
\subsubsection{基于特征交互的评分}
\label{sec:rs-feature-interaction}

线性模型$s=b+\bm w^\top\bm x$为每个特征赋予固定贡献，却不能单独表达“这类用户在这个页面更喜欢这件物品”。特征交互使一个特征的贡献依赖另一个特征。字段是应用定义的信息类别，如用户ID、物品ID和页面；一个字段可以包含许多类别特征。类别通常查表得到嵌入，连续数值则可以标准化后直接输入，或乘到其嵌入上。

下面的教学输入只激活三个特征，数值均为1，分别属于用户、候选与场景字段。为便于复算，嵌入暂取$\bm e_1=(1,0)$、$\bm e_2=(0,1)$、$\bm e_3=(1,1)$，偏置为$-2$，线性项之和为0.2。真实嵌入由标签学习，并没有把这两个维度预先命名为偏好。比较时需要看清模型交互的是标量、向量、字段对还是分支输出，以及汇总后还保留什么信息。

% Outline: 4.3.1.1.1
\paragraph{FM：二阶特征交互评分}
\label{sec:rs-fm}

线性模型为每个特征分配一个固定贡献，却无法单独表达“这类用户在这个页面更喜欢这件物品”的组合效应；若为每一对稀疏特征独立学习一个权重，$P$个类别就有约$P^2$个参数，训练中从未同时出现的组合也拿不到梯度。\term{因子分解机}（factorization machine, FM）的主要想法是让每个稀疏特征只维护一个潜在向量，用两向量内积充当该组合的二阶参数，把$P^2$次独立参数共享为$Pd$维嵌入，未观测组合也能借相邻向量取得分数\scite{rs-h05}

\textbf{模型结构。}第一步，对非零特征$a$查表得到嵌入$\bm e_a\in\mathbb R^d$，独热特征$x_a=1$，连续值可乘到嵌入上；第二步，线性项$b+\sum_a w_ax_a$给出各特征的固定贡献；第三步，二阶交互取所有非零特征对内积之和，形成完整评分
\begin{equation}
 s(\bm x)=b+\sum_a w_ax_a+
       \sum_{a<b}\langle\bm e_a,\bm e_b\rangle x_ax_b.
 \label{eq:rs-fm}
\end{equation}
用户特征与不同物品、场景交互时复用同一个$\bm e_a$，使未直接观测的组合也能借助共享参数取得分数。若输入只有用户和物品两个独热字段，交互部分就退化为两侧潜在因子的内积；加入页面、时间等字段后，则还会出现用户—页面、物品—页面等项。

第四步是平方改写以摊薄计算。设有$q$个非零特征，直接枚举特征对需$O(q^2d)$；展开平方时每个不同特征对出现两次，因此交互项可以改写为
\begin{equation}
 \frac12\sum_{r=1}^{d}
 \left[\left(\sum_a e_{a,r}x_a\right)^2
       -\sum_a e_{a,r}^2x_a^2\right].
 \label{eq:rs-fm-linear}
\end{equation}
减去自乘项后只剩交叉项，复杂度降为$O(qd)$。在教学输入中，三个内积为0、1、1，故$s=-2+0.2+2=0.2$，点击预测$\sigma(0.2)\approx0.5498$。

训练以点击二元交叉熵为目标并加嵌入正则，嵌入、线性项与偏置由同一条误差更新；服务时嵌入表可离线加载，线上只对激活特征算内积，不需要为新组合临时生成参数。若该样本点击为1，$\partial\ell/\partial s=p-y\approx-0.4502$，而
$\partial s/\partial\bm e_1=x_1\sum_{a\ne1}x_a\bm e_a=(1,2)$。
忽略正则、学习率取0.1时，$\bm e_1$一次更新为$(1.0450,0.0900)$；其他嵌入、线性项及偏置也由同一误差更新。

总类别数为$P$时，嵌入参数约为$Pd$，低秩共享使FM以远少于独立组合权重的参数表达二阶关系；但它只覆盖两两交互，三阶及以上作用留给后续模型，从未出现且没有辅助信息的特征仍缺少学习依据，低秩共享不能凭空解决所有冷启动。

% Outline: 4.3.1.1.2
\paragraph{FFM：字段感知交互评分}
\label{sec:rs-ffm}

FM要求一个特征在所有交互中复用同一个向量，可是同一用户与商品交互时需要的信息，可能不同于他与广告位交互时的信息：面向商品应突出偏好风格，面向广告位应突出出价与位置的适配。\term{字段感知因子分解机}（field-aware factorization machine, FFM）让一个特征拥有按对方字段区分的表示，把“和谁交互”纳入向量选择\scite{rs-h06}

\textbf{模型结构。}第一步仍为每个特征查表，但表项按对方字段复制：若$f(a)$为特征$a$所属字段，特征$a$面向字段$f(b)$的向量记为$\bm e_{a,f(b)}$。第二步，二阶项改为调用字段对各自的向量
\begin{equation}
 \sum_{a<b}
 \left\langle\bm e_{a,f(b)},\bm e_{b,f(a)}\right\rangle x_ax_b.
 \label{eq:rs-ffm}
\end{equation}
第三步，线性项、偏置与二阶和合成logit，点击输出和交叉熵可沿用FM，变化只在于每一对调用的参数。

例如用户面向物品的向量为$(1,0)$，物品面向用户的向量为$(0.2,0.5)$，贡献为0.2；用户面向页面的向量为$(0,1)$，页面面向用户的向量为$(0.3,0.7)$，贡献为0.7。FM要求用户在这两项中复用同一向量，FFM允许两种关系分开。一次点击误差只更新本样本实际使用的字段对向量及其他评分参数；服务时仍只计算激活特征对的内积，字段表项按字段数复制，存储随之增加。

若有$F$个字段、$P$个类别和$d$维表示，参数约为$PFd$；一般需要逐对计算，代价约为$O(q^2d)$，它失去了FM式~\eqref{eq:rs-fm-linear}所依赖的单一向量共享，因而不能直接套用平方改写。字段划分太细会使每个表项观测更少、参数更稀疏，正则与停止训练时机也更重要。额外自由度是否值得，应在相同输入、标签和资源下检验；字段定义是应用约定，更多参数不保证在所有稀疏数据上优于FM。

% Outline: 4.3.1.1.3
\paragraph{Wide \& Deep：组合评分}
\label{sec:rs-wide-deep}

FM只显式建模两两交互，更深的组合关系要靠网络间接拟合；而业务中已有一批被反复验证的明确组合（如“已安装应用A且候选为应用B”），直接给它一个可学习权重比让网络重新泛化更可靠。\term{Wide \& Deep}把显式线性分支与嵌入加深层分支相加为logit，并联合训练，让前者记住已见规则、后者在表示上泛化\scite{rs-r31}

\textbf{模型结构。}第一步，wide分支接收原始特征与人工显式交叉$\phi(\bm x)$，用线性权重直接打分；第二步，deep分支把类别嵌入及连续输入拼接后经MLP得到$\bm h_{\mathrm d}$；第三步，两路logit相加进入sigmoid
\begin{equation}
 p=\sigma\!\left(\bm w_{\mathrm w}^{\top}
       [\bm x;\phi(\bm x)]
       +\bm w_{\mathrm d}^{\top}\bm h_{\mathrm d}+b\right).
 \label{eq:rs-wide-deep}
\end{equation}
$\phi(\bm x)$是显式交叉特征，$\bm h_{\mathrm d}$是类别嵌入及连续输入经过MLP后的表示。显式分支可为“已安装应用A且候选为应用B”设置直接参数，深层分支则在应用、用户与情境表示上共享计算。原论文的应用是Google Play推荐。

假设教学线性分支logit为0.8，深层分支为$-0.3$，偏置为零，则$p=\sigma(0.5)\approx0.6225$；若点击为1，两个分支通过同一$p-y$得到梯度，分别调整交叉权重与嵌入网络。这里是联合训练的同一预测器，不是两个独立概率的平均；服务时wide分支按交叉字典查表加和，deep分支一次前向，两路共享同一条点击损失。

常见的“记忆与泛化”描述对应两种参数使用方式：显式组合可以直接记住已见关系，嵌入共享有机会支持相似但未见的组合。它没有保证新ID已有可靠表示；交叉字典膨胀、字段缺失及两分支时延也需要计入比较，记忆能力来自人工设计的交叉集，不等于冷启动已经解决。

% Outline: 4.3.1.1.4
\paragraph{DeepFM：低阶与高阶联合评分}
\label{sec:rs-deepfm}

Wide \& Deep的wide分支要人工选择哪些组合进入交叉字典，覆盖不全且维护成本高。\term{DeepFM}改用FM自动形成低阶交互，并让FM与深层分支共享原始特征嵌入，省掉显式交叉设计\scite{rs-r32}

\textbf{模型结构。}第一步，各类别特征查同一张嵌入表，FM分支按式~\eqref{eq:rs-fm}对嵌入内积求和产生低阶logit；第二步，deep分支把字段嵌入拼接成向量送入MLP，输出高阶logit；第三步，两logit相加进入sigmoid：
\[
 p=\sigma\!\left(s_{\mathrm{FM}}+
       \operatorname{MLP}([\bm e_1x_1;\ldots;\bm e_Fx_F])\right).
\]
同一个点击损失同时更新两条路径及共享嵌入。

在前述三字段例子中，FM分数为0.2。若深层输出0.3，则总概率为$\sigma(0.5)\approx0.6225$。相比Wide \& Deep，低阶交叉不再依赖单独指定每个组合，且两分支共享同一嵌入表，避免wide与deep各自维护一套表示。深层分支可以表示复杂非线性关系，但不能仅凭某一隐层就宣布识别了特定的三阶业务作用；其容量、数据需求和时延仍随网络配置改变，共享嵌入也不保证冷门ID得到有效训练。

% Outline: 4.3.1.1.5
\paragraph{DCN：有界阶特征交叉评分}
\label{sec:rs-dcn}

把嵌入直接拼接给MLP后，乘法关系只能靠网络逐层间接学习，既无显式结构也难以保证高阶组合被表示。\term{深度交叉网络}（deep \& cross network, DCN）显式让原始输入与当前表示逐层交叉，以紧凑参数表达有界阶交互，再与并行深层分支融合\scite{rs-r33}

\textbf{模型结构。}第一步，把原始输入$\bm z^{(0)}$同时送入交叉分支与深层分支；第二步，交叉层先把当前表示投影成一个标量$\bm w_\ell^\top\bm z^{(\ell)}$，再用它逐维缩放原始输入，加偏置与残差：
\begin{equation}
 \bm z^{(\ell+1)}
 =\bm z^{(0)}\bigl(\bm w_\ell^\top\bm z^{(\ell)}\bigr)
   +\bm b_\ell+\bm z^{(\ell)}.
 \label{eq:rs-dcn}
\end{equation}
第三步，将交叉表示与并行深层分支拼接，再接sigmoid预测头。若只看交叉分支，对原始连续输入的多项式最高阶每层增加一阶，但系数受共享结构约束，并非每个高阶单项式都有独立参数。

教学例取$\bm z^{(0)}=(1,2)$、$\bm w_0=(0.1,0.2)$、偏置为零。第一层标量为0.5，输出$(1.5,3)$；第二层仍取同一教学权重，标量为0.75，输出$(2.25,4.5)$。展开第二层可看到原始输入、二阶项与三阶项共同存在。将交叉表示与并行深层分支拼接，即可用点击交叉熵训练，梯度同时回到交叉层与深层分支。

若输入维度为$D$，原始交叉层的参数和乘加为$O(D)$；它比独立列举所有组合紧凑，也限制了可表示的交互权重。原论文属于ADKDD 2017工作坊。评价时应区分交叉层深度、深层分支容量和新增字段，不把“显式高阶”当作无条件优越性，也不把有界阶交叉误读为任意高阶组合都有独立参数。

% Outline: 4.3.1.1.6
\paragraph{DCN-V2：矩阵交叉与低秩评分}
\label{sec:rs-dcn-v2}

DCN每层用一个标量$\bm w_\ell^\top\bm z^{(\ell)}$缩放所有原始维度，所有输出维读到的是同一个组合强度，交互权重受到限制。\term{DCN-V2}把向量参数推广为矩阵，使不同输出维度读取不同的输入组合\scite{rs-r34}

\textbf{模型结构。}第一步，保留原始输入$\bm z^{(0)}$与残差连接；第二步，用矩阵把当前表示投影为同维向量，再与原始输入逐元素相乘：
\begin{equation}
 \bm z^{(\ell+1)}
 =\bm z^{(0)}\odot(\bm W_\ell\bm z^{(\ell)}+\bm b_\ell)
  +\bm z^{(\ell)}.
 \label{eq:rs-dcn-v2}
\end{equation}
注意偏置位于逐元素乘法内部，与式~\eqref{eq:rs-dcn}的原始写法不同。取$\bm z^{(0)}=(1,2)$及
$\bm W_0=\left(\begin{smallmatrix}0.1&0.2\\0.3&0.4\end{smallmatrix}\right)$，
偏置为零，则矩阵乘积为$(0.5,1.1)$，一层输出为$(1.5,4.2)$，第二维不再被同一个0.5缩放。第三步，交叉网络可与MLP并行后拼接，也可先交叉再堆叠MLP，最终用相同响应损失学习。

完整矩阵需要$O(D^2)$参数与计算；近似$\bm W=\bm U\bm V^\top$、秩为$r$时可降为$O(Dr)$。论文还研究低秩专家混合，在低维空间加入非线性和门控。这里的专家用于交叉变换，与MMoE按响应任务路由专家的职责不同。比较完整矩阵、低秩及混合方案，要共同记录误差、参数和候选批次时延。该工作正式发表于WWW 2021，2020年作者稿的日期不应替代会议年份。

% Outline: 4.3.1.1.7
\paragraph{xDeepFM：向量级显式交互评分}
\label{sec:rs-xdeepfm}

FM把一对嵌入的内积压成一个标量，整条向量各维度的组合关系被一次求和抹掉；DCN在标量输入维上交叉，也不区分嵌入维度。\term{xDeepFM}把一个字段的整条嵌入作为交互单位，由压缩交互网络（compressed interaction network, CIN）显式形成向量级交互\scite{rs-r35}

\textbf{模型结构。}令$\bm X^{(0)}\in\mathbb R^{F\times d}$为字段矩阵。第一步，上一层$\bm X^{(\ell-1)}$的每个向量与原始字段向量逐元素相乘，得到$H_{\ell-1}\times F$个乘积向量；第二步，每个输出通道对这些乘积向量做线性组合并压缩成一个新向量，第$h$个输出向量为
\begin{equation}
 \bm X^{(\ell)}_{h,:}
 =\sum_{a=1}^{H_{\ell-1}}\sum_{b=1}^{F}
   \bm W^{(\ell,h)}_{a,b}
   \bigl(\bm X^{(\ell-1)}_{a,:}\odot\bm X^{(0)}_{b,:}\bigr).
 \label{eq:rs-cin}
\end{equation}
同一字段对的权重跨嵌入维度共享，区别于对每个标量单独设交叉系数。第三步，各层输出沿嵌入维求和，汇总后与线性分支、深层分支共同产生点击logit。

在三字段例子中，若第一层仅给$(1,3)$和$(2,3)$两对权重1，其余为0，则输出为$(1,0)+(0,1)=(1,1)$，维度汇总为2。下一层再与原始字段相乘，可以进一步形成向量组合。该例只展示一个压缩通道，真实模型可学习多个通道并接CTR交叉熵。计算随$H_\ell H_{\ell-1}Fd$增加，压缩宽度是容量与代价的选择；它不是简单把FM叠得更深，也不消除其他分支的数据需求。

% Outline: 4.3.1.1.8
\paragraph{AutoInt：注意力特征交互评分}
\label{sec:rs-autoint}

FM、DCN、CIN的交叉权重对所有样本固定，可是同一候选在不同样本中应突出的字段并不相同。\term{AutoInt}把各类别、数值字段变成同维向量，让多头自注意力在字段之间按当前样本动态交换信息，再由预测头输出点击概率\scite{rs-r36}

\textbf{模型结构。}第一步，类别查嵌入、数值标准化后投影到同一维度；第二步，一个头内字段$a$以自身投影$\bm W_Q\bm e_a$为查询，对所有字段算兼容度并softmax：
\[
 \alpha_{a,b}=
 \frac{\exp(\langle\bm W_Q\bm e_a,\bm W_K\bm e_b\rangle)}
      {\sum_c\exp(\langle\bm W_Q\bm e_a,\bm W_K\bm e_c\rangle)},
 \qquad
 \widetilde{\bm e}_a=\sum_b\alpha_{a,b}\bm W_V\bm e_b.
\]
第三步，多头结果拼接，加上投影残差并经过激活，得到下一层字段表示；堆叠以后再接响应头并联合学习。这里写的是原模型的内积注意形式，不必强行添加其他Transformer变体的缩放。

教学例暂令投影为恒等，用户字段$(1,0)$对三个字段的logit为$(1,0,1)$，权重约为$(0.422,0.155,0.422)$，聚合为$(0.845,0.578)$。残差再加回用户字段后，原信息与其他字段信息共同进入后层。字段两两注意通常带来$O(F^2d)$量级的交互计算，投影另需预算。训练用点击交叉熵更新注意投影、残差与预测头；服务时字段集合在请求内固定，注意随当前样本重算。

AutoInt的注意对象是同一候选样本内的字段；DIN的注意对象是与候选匹配的历史行为。两者可以结合。注意权重有助于检查模型读取了什么，但不等于字段的因果作用，更不能据此推断用户真正作选择的心理过程。

% Outline: 4.3.1.1.9
\paragraph{RankMixer：异质特征混合评分}
\label{sec:rs-rankmixer}

扩大交互网络时，既要让字段交换信息，也要保留不同字段的独立变换能力；AutoInt的全局注意混合了表示，却没有为每个单元留下独立的非线性参数。\term{RankMixer}将用户、候选及已经加工好的历史表示分组为等维单元，通过多头重排混合交换子向量，再由每个单元独立的前馈网络变换\scite{rs-r39}

\textbf{模型结构。}第一步，把各组表示整理为等维单元；第二步，多头重排混合按头切分并交换各单元的子向量，只改变数据排列而不新增可学参数；第三步，交换后的子向量进入每个单元独立的前馈网络，这是容量的主要来源；层间保留残差及归一化，汇总后由任务头拟合点击等响应。

一个简化重排例子有两个四维单元$(a_1,a_2,a_3,a_4)$与$(b_1,b_2,b_3,b_4)$。按两头切分并交换后，可以形成$(a_1,a_2,b_1,b_2)$及$(a_3,a_4,b_3,b_4)$，各自进入独立MLP。于是一个单元的后续输出同时依赖两个原字段。实际模型按设定的头数和形状执行重排。

稀疏专家变体还需区分总参数与每个样本激活的参数。应在固定历史编码器、字段和标签下比较其与DCN-V2、AutoInt的交互代价。RankMixer仍可依赖独立序列模块先加工历史，并未单凭字段混合统一全部行为建模。文献按CIKM 2025会议记录使用，作者稿模板中的占位年份与DOI不构成发表证据。

表~\ref{tab:rs-feature-models}把前述计算放在共同坐标下。表中均可接点击概率头及交叉熵，区别在于交互和共享；这一共同输出不意味着原始输入、容量或训练预算天然相同。
\begin{table*}[htbp]
 \centering\small
 \begin{tabularx}{\linewidth}{@{}lXXXX@{}}
 \toprule
 方法 & 字段与交互 & 参数共享 & 汇总及输出 & 代价与适用条件\\
 \midrule
 FM & 稀疏特征的两两内积 & 每特征一个向量 & 二阶和加线性项 & 稀疏线性计算；低秩假设\\
 FFM & 对方字段参与选向量 & 按特征和对方字段 & 二阶和加线性项 & 更多表项；需字段对数据\\
 Wide \& Deep & 显式交叉及嵌入MLP & 组合权重与表示分工 & 两分支logit相加 & 交叉字典与深层计算\\
 DeepFM & FM及嵌入MLP & 两分支共用嵌入 & 两分支logit相加 & 深层增量成本；联合梯度\\
 DCN & 原输入乘当前投影 & 向量参数约束系数 & 交叉与深层融合 & 交叉层线性代价\\
 DCN-V2 & 矩阵或低秩交叉 & 矩阵、低秩或专家 & 并行或堆叠响应头 & 秩、专家与延迟取舍\\
 xDeepFM & CIN向量级乘积 & 跨维共享组合权重 & 层级汇总加其他分支 & 字段数和压缩宽度\\
 AutoInt & 字段间多头注意 & 共享投影和残差 & 保留字段后预测 & 字段两两读取\\
 RankMixer & 分组单元重排混合 & 单元前馈分别参数化 & 汇总后任务预测 & 激活计算与重排开销\\
 \bottomrule
 \end{tabularx}
 \caption{特征交互评分的共同接口。比较性能时还须固定输入信息、标签窗口及计算预算；不同共享限制没有脱离数据条件的统一优劣次序。}
 \label{tab:rs-feature-models}
\end{table*}

% Outline: 4.3.1.1.10
\paragraph{PNN：乘积层的特征交互评分}
\label{sec:rs-pnn}

FM在内积之后直接求和，乘积结果没有机会进入后续非线性网络；而把字段嵌入直接拼接给MLP时，乘性交互只能由网络隐式形成。\term{乘积神经网络}（product-based neural network, PNN）在字段嵌入与MLP之间加入乘积层，显式组合线性信号和乘积信号后预测响应\scite{rs-r46}

\textbf{模型结构。}第一步，各类别字段查表得嵌入；第二步，乘积层在字段对上计算交互：内积版本把字段对压成标量，外积版本保留不同潜在维度之间的组合；原论文用分解及汇总结构降低直接构造全部乘积的成本；第三步，乘积信号与线性嵌入信号拼接，送入多层网络输出响应。

对教学用户$\bm e_1=(1,0)$与对象$\bm e_2=(0,1)$，内积为0，而外积为
$\left(\begin{smallmatrix}0&1\\0&0\end{smallmatrix}\right)$。
内积丢弃了“第一维与第二维共同出现”的信息，外积仍保留它。若将三个字段的两两内积$(0,1,1)$送入教学隐单元，权重取$(1,0.5,-0.5)$且偏置为0.2，则ReLU输出0.2；以它直接作为logit得到概率0.5498，点击正例损失约0.598。实际PNN还并入线性信号并使用多层网络，用点击交叉熵联合训练乘积层与后续网络。

逐对内积约需$O(F^2d)$，逐对完整外积约需$O(F^2d^2)$；论文降本实现改变了参数化与汇总方式，比较时须指明版本。后续非线性响应头使它一般不能直接转成普通单向量点积索引召回，服务于给定候选评分时应按候选数计算预算。

% Outline: 4.3.1.1.11
\paragraph{NFM：交互池化的非线性评分}
\label{sec:rs-nfm}

FM把潜在维度也一并相加，交互被压成一个标量，后续网络便看不到不同维度的组合模式；PNN的外积又保留了完整矩阵、成本较高。\term{神经因子分解机}（neural factorization machine, NFM）把二阶交互先汇总成向量，再进行非线性变换，在保留维度的同时避免外积矩阵\scite{rs-r47}

\textbf{模型结构。}第一步，非零特征对逐元素相乘；第二步，双交互池化把所有对求和成一个向量
\begin{equation}
 \bm b(\bm x)=\sum_{a<b}x_a\bm e_a\odot x_b\bm e_b
 =\frac12\left[
 \left(\sum_a x_a\bm e_a\right)^{\odot2}
 -\sum_a(x_a\bm e_a)^{\odot2}\right].
 \label{eq:rs-nfm}
\end{equation}
右侧同样只需$O(qd)$。第三步，$\bm b$经非线性隐层，与偏置、线性项合成预测；无隐层且输出权重全为1时，恢复FM的二阶和。

三字段例子的交互向量为$(1,1)$。若教学输出权重为$(1,-1)$且没有隐层，则交互输出为0，总预测为$-1.8$；若此时回归标签为1，以半平方误差计损失为3.92。该人为权重用来说明两个维度可以承担不同作用，不代表论文训练结果。原论文以稀疏回归形式评估上下文预测与标签推荐；迁移CTR时，须再定义sigmoid头、点击标签及交叉熵。

NFM保留了潜在维度，却在池化时丢失了具体交互对的身份。若需要判断“哪一对字段更应参与”，就需要在汇总之前作选择，这一缺口由AFM补上。

% Outline: 4.3.1.1.12
\Needspace{90mm}
\paragraph{AFM：注意权重的二阶交互评分}
\label{sec:rs-afm}

NFM均匀汇总所有交互对，用户与对象的相关组合会与无关情境组合被同等对待。\term{注意因子分解机}（attentional factorization machine, AFM）为每个二阶交互向量$\bm b_{a,b}=x_ax_b(\bm e_a\odot\bm e_b)$计算权重，再加权汇总\scite{rs-r48}

\textbf{模型结构。}第一步，逐对形成交互向量$\bm b_{a,b}$；第二步，注意网络先产生未归一化分值
$a'_{a,b}=\bm h^\top\operatorname{ReLU}(\bm W\bm b_{a,b}+\bm c)$，
在本样本所有有效特征对上softmax得$a_{a,b}$；第三步，加权池化后经输出投影$\bm p^\top\sum_{a<b}a_{a,b}\bm b_{a,b}$，与偏置、线性项合成预测。注意网络、嵌入及预测头由任务损失共同学习。

教学三对向量为$(0,0),(1,0),(0,1)$。若注意权重为$(0.1,0.7,0.2)$，汇总为$(0.7,0.2)$；均匀权重则为$(1/3,1/3)$。以$\bm p=(1,-1)$投影，二者分别贡献0.5与0，说明样本相关权重能改变潜在维度的合成。

原论文采用稀疏回归，CTR迁移仍需明确更换输出及损失。与NFM相比，逐对注意增加了$O(q^2)$个交互对象的处理，不能直接沿用其全部线性化优势。高权重表示当前模型重视这一交互，不能证明该字段对造成了点击。

% Outline: 4.3.1.1.13
\paragraph{FiBiNET：字段选择与双线性交互评分}
\label{sec:rs-fibinet}

AFM在交互之后为组合赋权，却没有先判断字段本身在当前样本中是否重要；FM的内积也只是两向量逐维相乘，没有先在维度间混合。\term{FiBiNET}同时处理这两件事：交互前用压缩激励网络重加权字段，交互时用双线性变换代替纯内积\scite{rs-r49}

\textbf{模型结构。}第一步，压缩激励网络从当前字段嵌入提取摘要，生成字段权重$a_f$，得到重加权嵌入$\widetilde{\bm e}_f=a_f\bm e_f$；第二步，原始和重加权两组嵌入分别计算双线性交互；第三步，交互结果由浅层或深层预测头输出CTR。采用列向量记法，一对交互可写成$\bm b_{a,b}=(\bm W\bm e_a)\odot\bm e_b$。矩阵$\bm W$可以全局共享、按一个字段共享，或为每个字段对单独设置。它在逐元素相乘前混合潜在维度，区别于FM直接内积。

假设用户权重为0.5、场景权重为1.5，$\bm e_1=(1,0)$、$\bm e_3=(1,1)$，且$\bm W=\left(\begin{smallmatrix}1&0\\1&1\end{smallmatrix}\right)$，原始交互为$(1,1)$，重加权交互为$(0.75,0.75)$。两者拼接后仍可由预测头分别利用。字段门控、双线性矩阵与点击头通过交叉熵联合学习。

矩阵参数从$O(d^2)$增加到按字段的$O(Fd^2)$或按字段对的$O(F^2d^2)$，对应更细的关系和更高的数据需求。门控诊断、交互诊断及实际效果识别属于不同证据，不能互相替代。

% Outline: 4.3.1.1.14
\paragraph{Fi-GNN：字段图的交互评分}
\label{sec:rs-fi-gnn}

前述模型的交互多在一次池化中完成，字段关系只传递一轮；若希望每个字段经过多轮交换后再参与预测，可以把同一候选样本的字段构成图。\term{Fi-GNN}以字段为节点，以字段间的学习边权传播消息，再用GRU和初始状态残差更新节点\scite{rs-r50}这张图只有当前样本的用户、物品、页面等节点，并不是目录中的用户—物品交互图。

\textbf{模型结构。}第一步，节点初始化来自字段表示，边权由两端初始表示计算；第二步，每个节点具有输入和输出变换，边变换由发送端与接收端的矩阵组合，从而不必为每条边独立存一整套矩阵。以列向量约定写一轮消息：
\[
 \bm m_a=\sum_{b\ne a}\alpha_{b,a}
       \bm W^{\mathrm{in}}_a\bm W^{\mathrm{out}}_b\bm h_b,\qquad
 \bm h'_a=\operatorname{GRU}(\bm h_a,\bm m_a)+\bm h_a^{(0)}.
\]
第三步，若干轮传播后，节点预测与节点权重聚合成整体响应，用点击交叉熵训练。若教学变换为恒等，物品及场景传给用户的权重为0.6、0.4，则消息为$0.6(0,1)+0.4(1,1)=(0.4,1)$。GRU决定保留多少旧状态及新消息，初始残差帮助保留字段身份。

字段数为$F$时，全连接图仍有$O(F^2)$条边，每轮消息和状态更新均有成本；矩阵共享减少的是参数存储，并不使所有字段对消失。本节采用2019会议版机制，后续作者稿的修改应与它区分。节点、边的可视化说明模型的计算依赖，不构成因果解释。

% Outline: 4.3.1.1.15
\paragraph{DLRM：离散与连续特征的联合评分}
\label{sec:rs-dlrm}

前述模型多以类别嵌入为主，价格、时段统计等连续输入不宜总被当成类别；工业系统还需要在同一评分器里统一容纳离散与稠密两路。\term{深度学习推荐模型}（deep learning recommendation model, DLRM）让离散字段查嵌入表，连续输入经过底部MLP映射为同维稠密向量，再计算两两点积，并将点积及稠密向量交给顶部MLP\scite{rs-r51}

\textbf{模型结构。}第一步，离散字段各自查表得到嵌入；第二步，连续特征经底部MLP映射为与嵌入同维的稠密向量；第三步，所有向量做两两点积，保留每对内积而不像FM那样直接累加；第四步，点积与稠密向量拼接送入顶部MLP，输出logit。

例如用户、物品两条嵌入为$(1,0),(0,1)$，连续特征经底部MLP得到$(1,1)$。三个向量形成的非重复两两内积为$(0,1,1)$；顶部输入拼接为$(1,1,0,1,1)$，维度为$2+\binom32=5$。顶部MLP最后输出一个logit，sigmoid后以点击交叉熵训练，梯度同时到达稀疏表、底部及顶部网络。若稠密向量和$F$个稀疏向量都为$d$维，顶部输入维度为$d+\binom{F+1}{2}$。

FM直接累加各对交互，DLRM保留每对内积供顶部网络分别变换。显式操作虽然是二阶，经过非线性顶部后，整体函数不受二阶多项式限制。参数常集中于类别表，稠密计算集中于MLP与交互；两者并行实现的细节属于系统卷。改变采样或输入处理时仍需重新检查概率校准。

% Outline: 4.3.1.1.16
\paragraph{MaskNet：实例门控的特征评分}
\label{sec:rs-masknet}

FiBiNET对整条字段向量赋权，而同一用户换了候选，某些表示维度的重要程度也会变化，需要按实例逐维调节。\term{MaskNet}从当前原始嵌入生成逐维掩码，对嵌入或隐层作乘法缩放，并与前馈网络、层归一化组成MaskBlock\scite{rs-r52}

\textbf{模型结构。}第一步，压缩再扩展的网络从原始嵌入产生每个维度的掩码，通常不是0或1的硬选择；第二步，掩码与表示逐元素相乘；第三步，结果送入前馈网络并经层归一化，组成一个MaskBlock，每个块可从同一原始输入取得自己的掩码。串联MaskNet将一个块的输出交给下一个块，并联MaskNet让多个块处理共同输入，再拼接送入网络。

假设某层归一化表示为$(1,-1,0.5,-0.5)$。候选A对应的教学掩码为$(1,0.2,0.5,1)$，乘积为$(1,-0.2,0.25,-0.5)$；候选B掩码为$(0.2,1,1,0.5)$，乘积为$(0.2,-1,0.5,-0.25)$。即使用户相同，候选参与生成掩码后也会改变送入前馈层的表示。接预测头后即可按点击标签计算交叉熵并反向更新。

FiBiNET对整条字段向量赋权，MaskNet可对不同维度分别缩放。连续小权重不会自动跳过后续计算，也不能据此判定字段没有因果作用；门控处理与观测纠偏仍须分别检验。

% Outline: 4.3.1.1.17
\paragraph{Wukong：特征交互的容量扩展}
\label{sec:rs-wukong}

只扩大嵌入表，未必增加组合字段的计算能力；工业规模下需要沿稠密交互方向扩展容量，而不是单纯堆ID数量。\term{Wukong}以堆叠交互层扩展稠密计算，每层包含因子分解块与线性压缩块\scite{rs-r54}

\textbf{模型结构。}第一步，因子分解块从当前向量矩阵$\bm X$计算$\bm X\bm X^\top$等两两交互，将其展开、归一化、经MLP映射并重排为新向量；第二步，线性压缩块计算$\bm W_{\mathrm L}\bm X$，保留线性组合；第三步，两路拼接后与形状匹配的残差相加，再归一化：
\[
 \bm X^{(\ell+1)}=\operatorname{LN}\!\left(
 [\operatorname{FMB}(\bm X^{(\ell)});
  \operatorname{LCB}(\bm X^{(\ell)})]
  +\operatorname{Proj}(\bm X^{(\ell)})\right).
\]
这里方括号沿向量行拼接；形状相同时残差投影可以是恒等。堆叠多层后，最终响应头用点击等标签训练。

三字段例子的第一层内积矩阵为
$\left(\begin{smallmatrix}1&0&1\\0&1&1\\1&1&2\end{smallmatrix}\right)$。
因子分解块将这些关系变成新向量，压缩块则保留原字段的线性组合。第二层让两类向量继续相互作用，因此既能沿用低阶信息，也能组合前层已经形成的关系。由于包含MLP、归一化等非线性，交互阶数的直觉不应被误读为整个网络是一条固定阶多项式。

论文还利用低秩投影及乘法结合律减少显式内积矩阵的代价；输出向量数、层数、投影秩和MLP宽度共同控制预算。规模收益应限制在论文的数据及资源范围内，不能外推为无限增大必然改善。它扩展的是字段交互容量，字段—历史联合计算和跨阶段共享是另两个问题。

% Outline: 4.3.1.1.18
\paragraph{Hiformer：异质字段的注意力评分}
\label{sec:rs-hiformer}

AutoInt用同一套共享投影处理所有字段，可是应用ID、时段和地区虽都表示为向量，却不是语义相同的词元，比较时应读取的信息并不一样。\term{Hiformer}研究怎样为异质字段安排不同的信息投影\scite{rs-r55}

\textbf{模型结构。}基础异质注意层让字段$a$拥有自己的查询投影，字段$b$拥有自己的键、值投影，并采用字段专属前馈层。于是应用与时段比较时可以突出使用时间，应用与地区比较时可以突出地域可用性。完整Hiformer进一步使用复合投影：先拼接全部字段，再以一个分块矩阵生成各查询、键和值。例如
\begin{equation}
 [\bm k_1;\ldots;\bm k_F]=
 \widehat{\bm W}_{K}[\bm e_1;\ldots;\bm e_F],
 \label{eq:rs-hiformer}
\end{equation}
其中矩阵第$(a,b)$块允许输出单元$a$的键读取输入字段$b$。它不只是为单个字段各换一个矩阵，而是先生成含跨字段信息的复合单元，再进行注意交互。最后，加入的任务单元从交互层取得表示，由各自MLP头预测点击、安装或观看等目标。

一个标量教学投影可令应用键为$0.8e_{\mathrm{app}}+0.2e_{\mathrm{region}}$，时段键为$e_{\mathrm{hour}}$；相同应用在不同地区便能得到不同键。实际维度下，完整复合矩阵带来约$O(F^2d^2)$的投影成本，可通过低秩分解降低。末层只需任务单元输出时，可以裁剪不再用于预测的查询计算；这不等于前层从未计算过字段。本模型处理字段交互，并没有因此取得逐条长历史的顺序建模能力。

\subsubsection{基于历史行为的评分}
\label{sec:rs-history-scoring}

静态字段不能充分说明用户此刻在找什么。第~\ref{sec:rs-user-understanding}~节先把行为组织成候选无关的用户状态或向量；本节处理相反的接口：候选先进入历史计算，再输出该候选的响应分数。GRU4Rec、SASRec等模型先产生一份用户表示再与候选比较，属于上一章的表示接口；DIN以降的模型则等候选到来后，按候选激活、检索或压缩历史，得到的兴趣表示随候选变化，不能在请求前缓存成一份用户塔。

这些模型都以曝光响应损失训练，但监督位置随接口变化：短序列模型在给定候选上做激活或兴趣演化，长历史模型先为每个候选从完整历史选出相关子序列，再把候选与压缩后的历史共同送入响应头。若一份历史要比较几百个候选，独立用户编码只需计算一次，逐候选聚合则与候选数和读取的历史长度相乘，两种成本必须分开报告。最后一组STAN把候选相关匹配扩展到时间与空间条件，用于给定候选地点的响应评分。

\paragraph{DIN：候选相关的历史兴趣聚合}
\label{sec:rs-din}

在DIN之前，面向广告CTR的排序模型多采用“嵌入加多层感知机”的基础结构：各字段的稀疏取值先查嵌入表，用户历史这类多值字段再做一次与候选无关的求和或平均池化，压成一个固定长度的向量，最后与候选字段拼接送入多层感知机。这种池化带来两个问题。其一，无论评估哪个广告，读到的都是同一份历史表示，背包广告和相机广告无法各自调用相关行为；其二，用户兴趣是多样的，求和会把少量强相关事件淹没在大量无关事件中，若想让一个固定向量同时容纳所有兴趣，只能不断加大维度。\term{DIN}的主要想法是让候选参与历史读取：评估某个候选时，只突出历史中与它相关的行为，从而得到随候选变化的兴趣表示，再与其他字段共同预测CTR，即给定曝光条件下点击发生的概率\scite{rs-r12}

\textbf{模型结构。}DIN保留基础结构的嵌入层与预测网络，只替换用户历史的池化方式，图~\ref{fig:rs-din-activation}给出数据过程。第一步是嵌入：历史中第$k$条行为对应物品$i_k$，查表得到$\bm{v}_{i_k}$；当前候选$i$得到$\bm{v}_i$；用户画像、情境等其他字段各自查表后拼接为$\bm{x}_t$。第二步是局部激活：对每条行为，激活单元以行为向量、候选向量及二者的交互作为输入，经小型前馈网络输出一个标量权重。原文以两向量及其外积作为交互输入，常见实现改用差与逐元素积，可写为
\begin{equation}
  a_k(i)=g_\theta([\bm{v}_{i_k};\bm{v}_i;
            \bm{v}_{i_k}-\bm{v}_i;
            \bm{v}_{i_k}\odot\bm{v}_i]),\qquad
  \bm{h}(i)=\sum_{k=1}^{n}a_k(i)\bm{v}_{i_k}.
  \label{eq:rs-din-activation}
\end{equation}
第三步是加权求和，得到候选相关的兴趣向量$\bm{h}(i)$。第四步是预测：把$\bm{h}(i)$、$\bm{v}_i$与$\bm{x}_t$拼接，送入多层感知机，输出$\widehat p_i=\sigma(f_\theta(\cdot))$。

\begin{figure*}[htbp]
  \centering
  % DIN：给定候选对历史行为做局部激活（候选相关注意力），再加权求和，与其他字段拼接送入CTR头。
% 跨栏宽图（figure*），约150mm 宽。
\begin{tikzpicture}[x=1mm,y=1mm,
  font=\figurefont\footnotesize,
  dbox/.style={rectangle,draw=FigureInput!70,thick,fill=FigureInput!10,
    align=center,text=BookInk,rounded corners=1.4pt,minimum height=8mm,
    text width=22mm,inner sep=2pt},
  dhist/.style={rectangle,draw=FigureModel!70,fill=FigureModel!10,
    align=center,text=BookInk,rounded corners=1.4pt,minimum height=7.5mm,
    text width=20mm,inner sep=2pt},
  dact/.style={rectangle,draw=FigureLoss!75,fill=FigureLoss!14,
    align=center,text=BookInk,rounded corners=1.4pt,minimum height=8.5mm,
    text width=22mm,inner sep=2pt},
  dout/.style={rectangle,draw=FigureOutput!70,thick,fill=FigureOutput!12,
    align=center,text=BookInk,rounded corners=1.4pt,minimum height=10mm,
    text width=24mm,inner sep=2pt},
  dflow/.style={-{Stealth[length=1.8mm,width=1.2mm]},draw=FigureOutput!85,
    line width=.85pt},
  dfeed/.style={-{Stealth[length=1.8mm,width=1.2mm]},draw=FigureInput!85,
    line width=.9pt},
  dlab/.style={font=\figurefont\footnotesize,text=BookMuted,fill=white,inner sep=1pt}
]
  % 历史行为
  \node[dhist] (b1) at (14,18) {行为1 $i_1$\\背包};
  \node[dhist] (b2) at (14,9)  {行为2 $i_2$\\水壶};
  \node[dhist] (b3) at (14,0)  {行为3 $i_3$\\相机};
  \node[dhist] (b4) at (14,-9) {$\vdots$ 行为 $n$};
  \node[align=center,text=BookMuted] at (14,-15)
    {历史 $H_{u,t}$，嵌入 $\bm v_{i_k}$};

  % 当前候选与注入总线
  \node[dbox,fill=FigureInput!18] (cand) at (14,-25) {当前候选 $i$，$\bm v_i$};
  \coordinate (bus) at (37,-16);
  \draw[dfeed] (cand.north) -- (37,-21) -- (bus);
  \draw[dfeed] (bus) -- (37,22);
  \node[dlab,align=left,anchor=west] at (39,3) {候选 $\bm v_i$ 注入\\每个激活单元};

  % 行为 + 候选 -> 局部激活
  \node[dact] (a1) at (66,18) {局部激活\\$a_1(i)$};
  \node[dact] (a2) at (66,9)  {局部激活\\$a_2(i)$};
  \node[dact] (a3) at (66,0)  {局部激活\\$a_3(i)$};
  \node[dact] (a4) at (66,-9) {局部激活\\$a_n(i)$};
  \foreach \y in {18,9,0,-9}{
    \draw[dflow] (25,\y) -- (54,\y);
  }

  % 加权求和
  \node[dout,fill=FigureOutput!16,text width=30mm] (sum) at (100,5)
    {加权求和\\$\bm h(i)=\sum_k a_k(i)\bm v_{i_k}$};
  \foreach \y in {18,9,0,-9}{
    \draw[dflow] (78,\y) -- (88,\y) -- (88,5) -- (sum.west);
  }

  % 其他字段与预测头
  \node[dbox,fill=white,text width=26mm] (other) at (100,-13)
    {用户、情境等字段 $\bm x_t$};
  \draw[dflow] (other.north) -- (sum.south);
  \node[dout,fill=FigureOutput!20,text width=26mm] (mlp) at (134,5)
    {CTR 预测网络\\$\widehat p_i=\sigma(f_\theta)$};
  \draw[dflow] (sum.east) -- (mlp.west);
  \draw[dflow] (other.east) -- (127,-13) -- (127,-2) -- (mlp.south);
\end{tikzpicture}

  \caption{DIN的候选相关兴趣聚合。当前候选嵌入注入每个局部激活单元，与每条历史行为逐点比较得到权重$a_k(i)$，再对行为嵌入加权求和形成$\bm h(i)$；该向量与用户、情境字段拼接后送入CTR头。兴趣表示随候选变化，必须逐候选重算，不能预存为一份用户塔。}
  \label{fig:rs-din-activation}
\end{figure*}

这一结构有两处与普通注意力不同。权重$a_k(i)$不经softmax归一化，因此相关行为越多，$\bm{h}(i)$的幅度越大，强度信息得以保留；而归一化后的加权平均会让“三次浏览相机”和“一次浏览相机”给出相同方向和幅度。另外，求和与行为顺序无关，DIN不建模兴趣怎样随时间变化，这一缺口由DIEN补充。

用一个教学例子检查计算。假设登山包、水壶、相机的向量依次为$(1,0)$、$(0.8,0.2)$和$(0,1)$。对背包候选，权重为$(0.8,0.6,0.1)$，聚合得到$(1.28,0.22)$；对相机候选，权重为$(0.1,0.1,0.9)$，得到$(0.18,0.92)$。同一段历史对两个候选给出了方向不同的兴趣向量。若仅为演示取$f(\bm{h})=h_1+h_2-1$，两者CTR为$\sigma(0.50)$和$\sigma(0.10)$，约0.622和0.525；向量差异与最终概率排序由不同计算步骤决定，不能只看某个兴趣维度作结论。注意权重说明的是网络采用了哪些输入，不足以证明用户的真实动机。

训练以曝光点击的二元交叉熵为目标，嵌入表、激活单元与预测网络沿同一条损失联合更新。原文针对工业场景的稀疏特征补充了两项训练技术。小批量感知正则化只对本批出现过的稀疏ID施加与出现频率相关的惩罚，避免每步遍历整张亿级嵌入表；Dice激活按批统计量形成数据依赖的门，在两段斜率之间平滑切换，替代固定拐点的PReLU。这两项与候选相关聚合一起训练，但各自解决的是稀疏正则计算和非线性适配，不属于兴趣聚合本身。原工业实验用Adam、5,000样本批次和12维字段嵌入。

服务时稀疏ID只需查表，但激活单元依赖候选，必须在排序阶段对每个待评广告重新计算。因此DIN没有省掉“候选数$\times$历史长度”的在线匹配量，只是把这部分计算放在候选已经较少、信息最全的精排；兴趣向量也不能预先缓存成一份候选无关的用户塔。Dice在训练时用批统计量、服务时用累积统计量，服务端需要保留这组统计状态。输入候选集仍由召回决定，DIN不改变候选取得。

Alibaba实验用两周约20亿曝光训练、随后一天约1.4亿曝光测试，指标为按曝光数加权的用户AUC，与把全部曝光直接合并的AUC口径不同。在相同dropout与PReLU下，DIN从基础结构的0.5970升至0.6029；再加入小批量正则与Dice达到0.6083，消融可以分清历史激活和训练技术各自的贡献。2017年5至6月的显示广告A/B报告相对原线上模型最高10.0\% CTR、3.8\% RPM提升，并已服务主流量。该收益属于完整DIN及配套训练和服务方案，不能全部归于一组注意权重；公开电影评分二值化实验也不是实际广告点击测量\scite{rs-r12}

DIN的两个边界直接引出后续工作。求和不区分行为顺序，DIEN为此引入兴趣演化；逐候选扫描全部历史的代价随历史长度线性增长，实际系统通常只能读取最近一百多条行为，SIM为此先检索再精细匹配。

\paragraph{DIEN：候选相关的兴趣演化}
\label{sec:rs-dien}

DIN可以选择相关事件，却把历史当作无序候选集，未明确区分一串事件中的中间兴趣状态。\term{DIEN}先用GRU提取每一步状态，再让这些状态依据候选参与第二层演化。两层的分工是：第一层学习能解释随后行为的历史表示，第二层决定哪些变化与当前候选相关\scite{rs-r24}

\textbf{模型结构。}第一步，第一层GRU按时间读取行为嵌入，产出每步状态$\bm h_k$；第二步，为$\bm h_k$增加下一行为辅助任务，区分真实$i_{k+1}$与采样$\widetilde i_{k+1}$，让只受末端点击标签约束的早期状态也得到监督：
\begin{equation}
  \ell_{\mathrm{aux},k}
  =-\log D(\bm{h}_k,\bm{v}_{i_{k+1}})
   -\log[1-D(\bm{h}_k,\bm{v}_{\widetilde i_{k+1}})].
  \label{eq:rs-dien-aux}
\end{equation}
$D$是辅助判别网络。下一行为只用作监督，不能进入生成$\bm{h}_k$的前缀输入。第三步，第二层先计算候选与$\bm{h}_k$的注意权重$a_k(i)$，再通过注意更新门GRU（AUGRU）调节写入幅度，把候选相关性带进状态更新。若普通更新门为$\bm{z}_k$，则
\begin{equation}
  \widetilde{\bm{z}}_k=a_k(i)\bm{z}_k,\qquad
  \bm{g}_k=(\bm{1}-\widetilde{\bm{z}}_k)\odot\bm{g}_{k-1}
               +\widetilde{\bm{z}}_k\odot\widetilde{\bm{g}}_k.
  \label{eq:rs-dien-augru}
\end{equation}
AIGRU把注意权重乘到输入，弱输入仍可能触发状态变化；AGRU用标量注意力直接替换更新门；AUGRU保留各维门值，再整体缩放。假设$\bm{z}_k=(0.8,0.2)$、$a_k=0.25$，AUGRU得到$(0.2,0.05)$，而AGRU在两维都使用0.25，二者改变的计算不同。最终$\bm{g}_n$与用户、候选和情境送入CTR网络。

在登山与摄影交错历史中，相机候选可以降低与登山状态对应的写入，让摄影相关轨迹更多保留到末端。模型输出的是受点击及辅助任务约束的预测状态；称它为“兴趣演化”，并不意味着已经识别真实心理变化，更没有比较推荐行动怎样改变长期结果。与DIN相比，DIN用候选对整段历史加权求和、不建模状态随时间的传递，DIEN则在循环中按候选调制更新；后者代价是第二层必须逐候选运行，且多了辅助判别网络。

第一层GRU可与不同候选共享，但第二层AUGRU的门受到候选影响，必须逐候选运行。联合损失更新两层循环网络、注意、辅助判别器、嵌入及CTR头；辅助负例在Amazon实验从非目标物品随机采样，在工业实验取展示而未点击的广告，监督分布并不相同。推断移除辅助判别器及采样器，保留前缀的兴趣提取状态和候选相关演化路径，不能把未来辅助标签读回状态。

原论文Amazon Electronics与Books实验各重复5次，表4中Electronics普通双层GRU注意的AUC为0.7605，GRU加AUGRU为0.7640，加入辅助目标后为0.7792；Books相应为0.7890、0.7911、0.8453。工业CTR使用最近50条行为，表3的DIN、无辅助AUGRU、完整DIEN分别为0.6428、0.6493、0.6541，辅助任务增益明显小于公开数据。2018年6月7日至7月12日淘宝广告A/B相对BaseModel报告CTR提高20.7\%、eCPM提高17.1\%，同时PPC下降3.0\%；不能把广告主单次点击价格下降漏掉。部署还用核融合、批处理和蒸馏压缩，将报告的服务延迟从38.2毫秒降到6.6毫秒，这不是单靠AUGRU公式取得的延迟\scite{rs-r24}

\paragraph{DSIN：会话内兴趣与跨会话关系}
\label{sec:rs-dsin}

DIEN把整段历史视为一条循环链，默认相邻事件连续相关；行为按时间紧邻，并不一定来自同一个连续任务。上午集中选登山用品、晚上集中看摄影器材时，直接把全部事件看作一串等价位置可能掩盖会话结构。\term{DSIN}先按时间规则划分会话，在会话内提取兴趣，再建模会话之间的联系\scite{rs-r25}

\textbf{模型结构。}第一步，按时间规则把历史分为$m$个会话$H^{(1)},\ldots,H^{(m)}$，会话边界由预处理规则给定，不是注意力自动发现的真实意图切换；第二步，每个会话的行为嵌入加入区分会话、位置及特征维度的偏置编码，通过多头自注意力和池化形成会话内兴趣$\bm{s}_r$；第三步，双向LSTM继续读取请求前的全部会话向量，形成跨会话关系$\bm{g}_r$；第四步，候选对两组向量分别进行激活：
\begin{equation}
  \bm{h}_{\mathrm{session}}(i)=\sum_r\alpha_r(i)\bm{s}_r,\qquad
  \bm{h}_{\mathrm{relation}}(i)=\sum_r\beta_r(i)\bm{g}_r.
  \label{eq:rs-dsin-pooling}
\end{equation}
两组归一化权重由候选和相应会话表示的兼容程度决定。预测网络结合这两个输出及其他用户、物品字段，以曝光点击二元交叉熵训练全部可学习模块。

假设两个会话的向量分别为$(1,0)$与$(0,1)$，相机候选给它们的激活权重为0.2和0.8，则第一组聚合为$(0.2,0.8)$；第二组在双向LSTM之后计算，已经包含跨会话关系，不能直接复用这个数值。服务时同一候选同时读取两组聚合，而不是仅取最后一个会话。

双向计算仅覆盖预测时已知的历史会话。若会话划分把两个用途混在一起，或者因时间阈值把同一任务截断，模型建立的关系也会改变。因而应在相同事件上比较划分规则，并检验短会话、混合任务和边界附近的表现，而不能将“会话内部兴趣一致”当作全部数据都成立的事实。与DIEN的区别在于：DIEN在事件粒度上做候选相关演化，DSIN先把事件归并为会话单元、再在单元之间建模联系，粒度与组织依据不同。

会话内自注意参数跨会话共享，跨会话双向LSTM另有可学习门控，两组激活和最终MLP共同接受点击梯度。会话间隔阈值和最大行为数属于预处理选择，不通过这项梯度学习。服务先按相同边界规则编码已知会话，再对每个候选分别计算两组激活；候选无关的会话表示可复用，但边界追加或编码器更新后必须重新维护。

原文比较Alimama公开广告数据和Alibaba私有推荐数据，均取最近200条行为；前者约2,600万曝光、后者约60亿曝光，各用前7天训练、第8天测试。表1中偏置编码DSIN在广告数据的AUC为0.6375，对照DIEN为0.6343；推荐数据为0.6515与0.6473。移除会话关系层降至0.6365、0.6499，改成普通位置编码降至0.6357、0.6494。相同表中直接给DIN加RNN不如DIN，说明“先建立会话单元再建模关系”与“任何顺序层都有效”是不同结论。论文展示的是两个实际业务日志上的离线响应预测，并未在该实验表报告线上A/B收益\scite{rs-r25}

\paragraph{SIM：长历史的检索与精细匹配}
\label{sec:rs-sim}

DIN和DIEN在评估每个候选时都要逐条读取历史，计算量与历史长度成正比，实际系统通常只能读取不足150条的近期行为。可是与当前候选相关的证据可能出现在很久以前：用户半年前浏览过一批镜头，今天的相机广告恰好与之相关，而最近一百多条行为里全是日用品。MIMN尝试在行为到来时把历史压缩进固定容量的记忆，从而把可用历史延长到约1000条；但写入记忆时不知道未来会评估哪个候选，长度再增加十倍后，压缩会把候选需要的细节混掉。\term{SIM}的主要想法是保留原始长历史，把“从中找出相关行为”与“精细建模相关行为”拆成两个级联阶段：先用便宜的方法按候选检索出几百条相关行为，再只对这一小段做DIN式的候选相关注意力。这里检索的是用户已经发生的行为，不是待推荐的目录物品；后者是第~\ref{chap:rs-retrieval}~章召回的任务\scite{rs-r20}

\textbf{模型结构。}SIM由通用检索单元（general search unit, GSU）和精细匹配单元（exact search unit, ESU）级联而成，图~\ref{fig:rs-sim-two-stage}给出数据过程。设用户长期行为序列为$b_1,\ldots,b_n$，$n$可达上万。GSU为每条行为计算一个与候选$i$的相关分数$r_k(i)$，取分数最高的$K$条组成子序列$S_i$，$K$通常为几百；ESU以$S_i$为输入，用注意力形成长期兴趣向量；该向量与短期行为兴趣及其他字段拼接，送入多层感知机预测CTR。两级的分工依据计算预算确定：GSU面对全部历史，只能使用每条成本极低、可以预先建立索引的相关性；ESU只面对$K$条行为，可以使用DIN、DIEN这类较昂贵的结构。GSU的作用还不只是节省计算——长历史中大量与候选无关的行为本身就是噪声，原文的消融显示，即使只对筛选后的行为做平均池化，效果也明显好于直接平均全部长期行为。

\begin{figure*}[htbp]
  \centering
  % SIM：候选先在长历史上做通用检索 GSU 得到子序列，再由 ESU 精细注意力聚合，接入CTR。
% 跨栏宽图（figure*），内容宽约155mm；上行处理链，下方预维护与短期特征。
\begin{tikzpicture}[x=1mm,y=1mm,
  font=\figurefont\footnotesize,
  sstore/.style={rectangle,draw=FigureModel!70,fill=FigureModel!10,align=center,
    text=BookInk,rounded corners=1.4pt,minimum height=11mm,
    text width=24mm,inner sep=2pt},
  sgsu/.style={rectangle,draw=FigureLoss!80,thick,fill=FigureLoss!14,align=center,
    text=BookInk,rounded corners=1.4pt,minimum height=14mm,
    text width=27mm,inner sep=2pt},
  sesu/.style={rectangle,draw=FigureInput!75,thick,fill=FigureInput!12,align=center,
    text=BookInk,rounded corners=1.4pt,minimum height=14mm,
    text width=31mm,inner sep=2pt},
  sout/.style={rectangle,draw=FigureOutput!75,thick,fill=FigureOutput!14,align=center,
    text=BookInk,rounded corners=1.4pt,minimum height=12mm,
    text width=18mm,inner sep=2pt},
  slow/.style={rectangle,draw=FigureModel!60,fill=FigureModel!7,align=center,
    text=BookInk,rounded corners=1.4pt,minimum height=9mm,
    text width=27mm,inner sep=2pt},
  sflow/.style={-{Stealth[length=1.8mm,width=1.2mm]},draw=FigureOutput!85,
    line width=.9pt},
  sfeed/.style={-{Stealth[length=1.8mm,width=1.2mm]},draw=FigureModel!80,
    line width=.85pt,dashed},
  slab/.style={font=\figurefont\scriptsize,text=BookMuted,fill=white,inner sep=1pt}
]
  % 上行处理链
  \node[sstore,fill=FigureInput!16,draw=FigureInput!70] (cand) at (18,12)
    {当前候选 $i$\\$\bm v_i$ 与时间};
  \node[sgsu] (gsu) at (51,12)
    {通用检索 GSU\\硬检索：类别倒排\\软检索：内积};
  \node[sstore,fill=FigureLoss!8,draw=FigureLoss!70] (sub) at (85,12)
    {相关子序列 $S_i$\\Top-$K$ 行为与时间};
  \node[sesu] (esu) at (122,12)
    {精细匹配 ESU\\多头候选相关注意力\\仅在 $S_i$ 上聚合};
  \node[sout] (ctr) at (154,12) {CTR 网络\\点击概率};
  \draw[sflow] (cand.east) -- (gsu.west);
  \draw[sflow] (gsu.east) -- (sub.west);
  \draw[sflow] (sub.east) -- (esu.west);
  \draw[sflow] (esu.east) -- (ctr.west);
  \draw[sflow] (cand.north) -- (18,24) -- (122,24) -- (esu.north);

  % GSU 下方：长历史与索引预维护
  \node[slow] (hist) at (51,-8) {长期行为序列};
  \node[slow] (index) at (51,-19) {行为索引／嵌入（倒排或ANN）};
  \draw[sfeed] (index.north) -- (hist.south);
  \draw[sflow] (hist.north) -- (gsu.south);
  \draw[sflow] (index.east) -- (65,-19) -- (65,5) -- (gsu.south);
  \node[slab,align=center] at (51,-26) {行为到来时预先维护；原系统覆盖约180天、最多$5.4\times10^4$条};

  % ESU 下方：短期特征
  \node[slow] (short) at (122,-8) {短期兴趣及其他特征};
  \draw[sflow] (short.north) -- (esu.south);
  \draw[sflow] (short.east) -- (146,-8) -- (146,6) -- (ctr.south);

  % 关键标注
  \node[slab,align=center,text=FigureLoss!85] at (85,20)
    {截断边界：未入选行为不进 ESU，漏检不可恢复};
  \node[slab,align=center] at (51,2)
    {每候选计算从扫描 $n$ 条降为选 $K\ll n$ 条};
\end{tikzpicture}

  \caption{SIM的两阶段长历史匹配。GSU依据候选从预先维护的长期行为索引中选出相关子序列$S_i$，ESU只在这$K\ll n$条行为上做多头候选相关注意力，再与短期兴趣共同送入CTR。候选同时驱动GSU与ESU；未被GSU选中的行为不进入ESU，截断造成的漏检无法靠后阶段恢复。}
  \label{fig:rs-sim-two-stage}
\end{figure*}

GSU的逻辑是一次以候选为查询、在用户自己的历史上执行的检索。原文给出两种相关分数，硬检索与软检索：
\begin{equation}
  r_k(i)=
  \begin{cases}
    \mathbb{I}\{c(b_k)=c(i)\}, & \text{硬检索},\\
    \langle\bm{W}_{\mathrm b}\bm{e}_k,\ \bm{W}_{\mathrm a}\bm{v}_i\rangle, & \text{软检索},
  \end{cases}
  \label{eq:rs-sim-gsu}
\end{equation}
其中$c(\cdot)$为物品类目，$\bm{e}_k$为第$k$条行为的嵌入，$\bm{W}_{\mathrm b},\bm{W}_{\mathrm a}$为投影矩阵。

硬检索只保留与候选同类目的行为。它没有可学习参数，相关性完全由类目体系决定，因而可以离线为每个用户建立“用户—类目—行为列表”的两级索引，线上只需一次按类目的查表。它适用于类目划分细致、同类目物品大体相似的场景，例如原文所在的电商广告。其局限是跨类目但语义相关的行为会被漏掉，类目过粗时又会混入大量无关行为；同一类目下行为过多时还要另定截断规则。

软检索用学习得到的嵌入计算候选与行为的内积，再借近似最大内积检索（maximum inner product search, MIPS）在亚线性时间内取Top-$K$。它能发现跨类目的相关行为，适用于类目体系粗糙或物品关系主要体现在交互中的场景。代价有三项：嵌入需要训练，线上要维护向量索引并执行近邻检索；嵌入更新后索引也要同步；以及长期行为与近期行为分布不同，直接沿用短期兴趣模型的嵌入可能误导检索。为此原文为软检索单独设一个辅助CTR任务：把全部长期行为按$r_k(i)$加权求和为$\bm{u}_{\mathrm r}(i)=\sum_k r_k(i)\bm{e}_k$，与候选向量拼接后预测点击，用这条损失训练检索所用的嵌入；历史过长时可随机抽取子序列计算该项。

两种检索的取舍见表~\ref{tab:rs-sim-gsu-comparison}。原文在工业数据上发现，软检索选出的Top-$K$与硬检索非常接近，硬检索保留的行为覆盖了软检索结果的约75\%；软检索的离线AUC仅略高，却需要在线近邻检索和更多存储，因此部署采用硬检索。这个结论依赖“同类目物品多数相似”的数据特点，换到类目体系不同的场景需要重新比较。

\begin{table}[htbp]
  \centering
  \small
  \begin{tabularx}{\linewidth}{@{}lXX@{}}
    \toprule
    维度 & 硬检索 & 软检索 \\
    \midrule
    相关分数 & 与候选类目是否相同 & 投影嵌入的内积 \\
    可学习参数 & 无 & 嵌入与投影矩阵 \\
    训练 & 不需要，GSU损失权重为0 & 辅助CTR任务，与ESU共享嵌入 \\
    线上实现 & 离线两级索引，按类目查表 & 向量索引与近似最大内积检索 \\
    适用场景 & 类目细致、同类目物品相似 & 类目粗糙、跨类目关联重要 \\
    主要局限 & 漏掉跨类目相关行为 & 计算存储更高，索引需随嵌入同步 \\
    \bottomrule
  \end{tabularx}
  \caption{SIM通用检索单元的两种实现。两者都只决定哪些行为进入ESU，不改变ESU的精细匹配方式。}
  \label{tab:rs-sim-gsu-comparison}
\end{table}

ESU的逻辑是在已经筛出的$S_i$上，像DIN那样让候选决定每条行为的贡献，但增加了两项长历史特有的设计。第一，所选行为可能横跨数月，相隔久远的行为贡献不同，因此把行为与请求之间的时间间隔$\Delta_k$离散化后嵌入，与行为嵌入拼接为$\bm{z}_k=[\bm{e}_k;\bm{e}^{\mathrm t}_{\Delta_k}]$。第二，长期兴趣往往有多个方面，因此使用多头注意力：
\begin{equation}
  \alpha_{r,k}(i)=\operatorname*{softmax}_{k\in S_i}
     \langle\bm{W}_{\mathrm b,r}\bm{z}_k,\ \bm{W}_{\mathrm a,r}\bm{v}_i\rangle,\qquad
  \bm{h}(i)=\mathop{\Vert}_{r=1}^{q}\sum_{k\in S_i}\alpha_{r,k}(i)\bm{z}_k,
  \label{eq:rs-sim-interest}
\end{equation}
其中$q$为注意头数，$\Vert$表示拼接。与DIN不同，这里每个头的权重在子序列上作softmax归一化。由于输入已缩短到几百条，ESU也可以直接换成DIEN等更复杂的结构。

用一个例子比较两种检索的结果。设候选为一台相机，类目为“相机”；历史六条依次为登山包、水壶、镜头、帐篷、相机、摄影背包，类目分别为箱包、户外、相机、户外、相机、箱包。硬检索取出第3、5条，摄影背包因为属于箱包类目被漏掉；软检索若学到摄影背包与相机常被一起浏览，就能把它选入，但也可能因外形相近而误选登山包。无论哪种检索，未入选的行为都不会进入ESU，后阶段无法补回。

两个单元在同一个模型中联合训练，目标为$\mathcal{L}=\alpha\mathcal{L}_{\mathrm{GSU}}+\beta\mathcal{L}_{\mathrm{ESU}}$，两项都是曝光点击交叉熵。软检索时$\alpha=\beta=1$，GSU与ESU共享嵌入，每次检索使用最新参数得到Top-$K$；硬检索没有参数，取$\alpha=0$。需要注意，可学习的嵌入接受点击监督，并不意味着离散Top-$K$已经成为可微操作：ESU的梯度只流经被选中的行为，漏检行为得不到精细监督。

服务端的难点在于存储与时延。原系统高峰每秒服务百万以上用户，每个请求要为数百个候选预估CTR，延迟需低于30毫秒。部署方案为每个用户离线建立两级的用户行为树：第一级键为用户，第二级键为类目，值为该类目下的行为列表，整体为分布式存储，约22 TB，按天更新。线上GSU以候选类目查表，数百个候选通常只涉及不到20个类目，每个类目的子序列再截断到200条，单次请求的数据量因而可控；ESU的多头注意力另做了核融合优化。在这一方案下，SIM读取上万条历史，延迟仅比读取截断至1000条的MIMN增加约5毫秒。

工业实验使用49天曝光训练、随后一天测试，长期行为覆盖此前180天，短期行为覆盖此前14天，最长达54,000条，超过30\%的样本有上万条历史。表4中DIEN的AUC为0.6452，MIMN为0.6541，硬检索SIM为0.6604，软检索SIM为0.6625，硬检索加时间信息为0.6624。公开Taobao数据的消融中，仅软检索加平均聚合为0.9357，加入精细第二阶段为0.9416，再加时间信息为0.9501，说明筛选、精细聚合和时间各自改变了不同计算。2020年1月7日至2月7日的线上A/B相对MIMN报告CTR提高7.1\%、RPM提高4.4\%，对应的是硬检索加时间信息的完整部署路径\scite{rs-r20}

SIM与DIN的关系可以概括为：两者都让候选决定历史中哪些行为重要，区别在历史入口。DIN对全部行为逐条激活，每个候选的精细计算为$O(n)$；SIM先用GSU把范围缩到$K\ll n$条，精细计算降为$O(K)$，代价是引入一次不可恢复的截断。这次截断还带来一个新问题：GSU用类目或独立嵌入衡量相关性，ESU用注意力衡量相关性，两者标准不同，GSU可能选入ESU认为无关的行为，而漏掉ESU真正看重的行为。TWIN针对的正是这种两阶段相关性不一致。

\paragraph{TWIN：两阶段历史选择的相关性一致性}
\label{sec:rs-twin}

SIM两阶段还有一个接口问题：粗阶段按类别或独立内积选出的事件，可能与精细阶段注意力实际重视的事件不同，两套相关性标准不一致。\term{TWIN}让两阶段采用相同的候选—行为相关性计算，通过拆分特征降低长历史扫描成本。这样做的目标是让有限的精细计算留给同一度量认为相关的事件\scite{rs-r28}

\textbf{模型结构。}第一步，把历史特征拆成两类：物品自身特征的投影可以跨用户复用，点击时间、观看情况等用户—物品交叉特征则随事件变化；第二步，自身特征投影后缓存，交叉特征各自压成低维偏置，两阶段共用同一个相关性分数。省略注意头编号后，一条历史的相关性可概括为
\begin{equation}
  a_k(i)=
  \frac{\langle\bm{W}^{\mathrm{q}}\bm{v}_i,
                 \bm{W}^{\mathrm{h}}\bm{f}_{i_k}\rangle}{\sqrt d}
       +\bm{\beta}^{\top}\bm{c}_{u,k},
  \label{eq:rs-twin-score}
\end{equation}
其中$\bm{f}_{i_k}$为物品自身特征，$\bm{c}_{u,k}$为压缩后的事件交叉特征，$\bm{\beta}$为可学习参数。第三步，粗阶段据此选择历史，精细阶段在选集内对相同相关性作softmax，再聚合完整值投影；多个头的结果与近期行为等分支共同预测CTR。

假设相机候选对镜头、相机包和水壶的精细相关性为3、2、0，另一套粗度量却给出1、0、2。只保留两条时，后者漏掉相机包。统一度量保留镜头和相机包，精细权重为$e^3/(e^3+e^2)\approx0.731$和0.269。例子说明一致性解决了什么，也说明它没有证明第三条完全无用：Top-$K$仍然截断了尾部证据。与SIM相比，SIM的GSU与ESU各用一套相关性（类目或独立嵌入内积 vs 注意力），TWIN把两者统一为同一公式，代价是交叉特征仍要随事件压缩、缓存需要定期刷新。

点击监督训练共享相关性与预测网络，服务端复用其参数和物品投影。参数共享、参数同步及投影缓存刷新是不同要求：两阶段公式相同，缓存过时仍可让实际得分不同。相关性本身不准时，一致地选错也无济于事；因此应同时检查选择重合程度、最终预测和缓存时效，而不是只报告一致率。

多头粗阶段分别产生事件次序，原配置遍历四头结果并去重，直到取得100条行为；精细阶段按同头相关性聚合所选完整值。嵌入用AdaGrad更新，稠密网络用Adam更新，服务每15分钟刷新缓存的物品投影。交叉特征保留为低维偏置，使候选无关物品投影能缓存；删除这些事件字段和压缩它们不是同一加速办法。

Kuaishou原实验每天约462亿日志，历史平均14,500条、截断上限100,000条；主比较粗阶段读取最近10,000条、细阶段100条，连续5天各用23小时训练、最后1小时测试。这里GAUC是按用户样本数加权的用户内AUC，衡量同用户正负曝光的次序；表4中TWIN的AUC为0.7962、GAUC为0.7336，SIM Soft为0.7939、0.7299。表中的SDIM改用目标注意作为ESU，不能直接当作SDIM原桶聚合方案的比较。以各模型自身ESU定义的Top-100作检索参照，TWIN命中94条而非理论100条，差距与15分钟缓存刷新有关。参数不共享和删除交叉偏置的消融均降低效果。三入口线上A/B相对SIM Soft的观看时长提高1.374\%至2.778\%，并报告已部署主流量；这里测的是观看总时长，不是CTR或长期满意度\scite{rs-r28}

\paragraph{TransAct：即时行为与长期表示的组合}
\label{sec:rs-transact}

批量更新的长期用户向量按天或小时刷新，用户刚刚收藏摄影内容时，这一表示可能要稍后才会反映这一变化。\term{TransAct}把即时动作作为额外输入，与批量长期表示共同参与排序。它并不要求先重训整个长期模型，而是直接让当前请求读取近期可用证据\scite{rs-r14}

\textbf{模型结构。}第一步，每条近期事件包含物品内容向量、动作类型和时间戳，把动作嵌入与内容向量结合，再把候选向量拼接到每条行为上，形成输入$\bm{x}_k(i)=[\bm{v}_{i_k};\bm{e}_{b_k};\bm{v}_i]$。这种早融合不同于只把候选追加成序列末尾的一个位置：前者从每个输入位置就带有候选信息，后者依赖后续注意力交换。原工作比较后选择了逐行为拼接。第二步，Transformer编码按时间倒序排列的近期输入；第三步，取最近$K$个输出，并拼接全序列的逐维最大池化：
\begin{equation}
  \bm{h}_{\mathrm{recent}}(i)=
    [\bm{o}_1;\ldots;\bm{o}_K;
      \operatorname{maxpool}(\bm{o}_1,\ldots,\bm{o}_n)].
  \label{eq:rs-transact-compress}
\end{equation}
第四步，所得紧凑表示与长期表示及其他字段一起进入排序预测网络，用不同动作标签的加权预测损失训练。原配置未采用位置编码；时间戳用于排序和时间窗口掩码。因此，不能仅凭“使用Transformer”就宣称网络完整保留了全部位置关系。

训练随机抽取一个时间窗口，屏蔽请求前该窗口内的近期事件，促使模型也利用较远证据；推断不使用这个随机时间窗口掩码。它不同于防止未来泄漏的因果掩码，二者目的和部署方式都不同。假设近期两个输出为$(0.2,0.8)$与$(0.7,0.3)$、$K=1$，压缩结果为$(0.2,0.8,0.7,0.8)$，近期重点和逐维最大值都进入下游。

刚到达的收藏能改变这个分支，即使长期向量仍是昨日版本。但事件是否已被服务端读到、长期向量何时刷新、预测参数多久更新，依然是三条不同时间线。近期窗口也会遗漏更早事件；即时性改善需要在同一请求和标签窗口下验证，不能直接推导长期收益。与只靠批量长期向量的方案相比，TransAct用一条实时分支补时效，而不是用实时分支替代长期表示。

原配置读取100条近期行为，动作嵌入32维，使用两层Transformer并保留最近10个输出及最大池化。内容与动作编码、早融合、压缩以及下游多头共同接受响应梯度；既有批量用户向量作为输入，并不在每次请求重新训练。训练对负例下采样，离线评价则使用随机打乱位置的小流量真实日志，不能把训练正负比例当作总体比例。服务另由实时事件处理去重并对齐时间，模型才能读到这些动作。

Pinterest实验用两周约30亿实例训练、第三周评价。其HIT@3按同请求的推荐块计算，并按用户汇总，不是单个后继的HR@3。表1相对不含实时序列的Pinnability，TransAct的HIT@3/repin提高9.40\%、HIT@3/hide降低14.86\%；BST全动作版本反而增加hide，说明负动作类型必须被正确表达。表2去掉TransAct或PinnerFormer都降低效果，两者互补；保留所有输出也未同时改善保存与隐藏两项指标。每组1.5\%用户的线上测试报告Homefeed保存量提高11\%，但固定模型收益随时间衰减，重训才能维持。原文也报告已用于多个Pinterest推荐入口，这不意味着一次参数训练即可长期使用不断刷新的行为\scite{rs-r14}

\paragraph{HSTU：异质行为的序列转导}
\label{sec:rs-hstu}

前述方案把历史当作物品序列，动作类型只作辅助特征；推荐输入既有物品，也有点击、收藏和跳过等动作，不同任务又反复读取同一历史。\term{HSTU}把这些信息组织为异质序列，研究怎样共享多个合法位置的序列计算。序列转导在这里表示从已知事件序列预测后续物品或候选动作，输出对象由任务定义，而不是自由生成一个商品名称\scite{rs-r07}

\textbf{模型结构。}第一步，把物品标识与动作类型、时间等组织为输入矩阵$\bm{X}\in\mathbb{R}^{n\times d}$，合法位置由可见性掩码限定；第二步，一个HSTU层用逐位置投影同时生成查询、键、值及门控张量，并以位置和时间偏置调整聚合。省略多头重排后可写成
\begin{align}
  (\bm{U},\bm{V},\bm{Q},\bm{K})
    &=\operatorname{Split}\bigl(\operatorname{SiLU}(f_1(\bm{X}))\bigr),\\
  \bm{A}&=\operatorname{SiLU}(\bm{Q}\bm{K}^{\top}+\bm{B}),\\
  \bm{Y}&=f_2\bigl(\operatorname{Norm}(\bm{A}\bm{V})\odot\bm{U}\bigr).
  \label{eq:rs-hstu}
\end{align}
可见性掩码排除不合法的事件关系，层间使用残差连接。$\bm{B}$编码相对位置与时间，$f_1,f_2$为逐位置映射。这里逐对应用SiLU，没有将整行权重softmax归一化。第三步，同一序列编码接两种任务头：召回头在合法前缀位置预测下一物品，排序头把待判断物品放在相应输入位置、用当时尚未知晓的动作作标签。

这一差别影响证据数量的表达：若两个完全相同的值都获得相同softmax权重，它们的加权平均与只有一个值时相同；未作行归一化的聚合则可随重复证据增加幅度。HSTU随后仍有归一化、门控和投影，因此实际输出怎样保留强度要由整层计算判断，不能把未归一化直接解释为次数计数器。

召回训练在合法前缀位置预测下一物品，用户状态与物品输出表示构成目录分布，实际计算可采用采样目标。排序训练则把待判断物品放在相应输入位置，用当时尚未知晓的动作作标签，接入响应预测头。历史中的既有动作可作为输入，当前候选的待预测动作不能提前写入。一次浏览后收藏与浏览后跳过可以共享物品信息，却必须形成不同的动作监督。

例如，历史为“背包浏览、水壶收藏、相机跳过”，召回头回答下一件可能交互的物品，排序头回答给定新背包曝光后的某种动作概率。两者可以复用序列编码结构，但候选是否输入、损失位置和输出头都不同。推断时能复用多少历史计算，还取决于候选之间和候选对历史的可见性。稀疏训练、随机长度和M-FALCON等服务设计属于系统实现；论文在特定规模上的扩展结果不能被省去条件地推广到所有数据和预算。

召回在批次中组织多个合法前缀目标，以采样softmax更新序列层和物品输出表；排序在候选位置接动作头，用观测响应监督，同时掩去该位置尚未知晓的动作。原工业训练以流式单遍为主，通过随机长度减少参与训练的位置，不能把被省略位置当作零标签。服务保留每层历史键值和合法位置状态，以M-FALCON将多候选评分分块复用历史计算；只有掩码阻止候选回写历史、参数与历史版本一致时，这种复用才成立。

原论文分开报告公共序列数据与工业流式数据。公共数据表4在MovieLens-1M全目录评价中，HSTU的HR@10为0.3097、NDCG@10为0.1720，改进训练配方的SASRec为0.2853、0.1603；不是与2018年采样候选数字直接比较。工业表5去掉相对偏置或改回softmax均损害部分任务，说明层结构和输入组织各有作用。表6加入一条新召回路与替换原路分别测得不同线上收益；表7排序GR相对多年优化的DLRM在E-Task、C-Task报告12.4\%、4.4\%提升，并已在Meta多个应用部署。不过公开材料仅用这两个任务代称，未给出足够业务指标定义，不能将12.4\%擅自写成CTR、收入或留存。H100上编码器最高15.2倍训练、5.6倍推断加速还限定了长度、宽度、稀疏策略和算子实现，并非端到端服务的统一倍数\scite{rs-r07}

\paragraph{ETA：候选相关的高效长历史选择}
\label{sec:rs-eta}

SIM硬检索依赖类目体系、软检索需要维护向量近邻索引；长历史检索还可以使用更便宜的二进制比较，但检索向量仍应反映点击任务。\term{ETA}用CTR网络中的物品嵌入生成SimHash指纹，通过候选与行为指纹的汉明距离选择相关历史，再对选中事件作目标注意力，与短期分支共同预测点击\scite{rs-r76}

\textbf{模型结构。}第一步，取CTR网络中训练得到的物品嵌入，经固定随机投影方向生成二进制指纹；第二步，以候选与行为指纹的汉明距离近似向量方向差异，选出距离最小的$K$条；第三步，对选中事件作目标注意力，与短期分支共同预测点击。给定随机投影方向$\bm{r}_q$，第$q$位指纹为
\begin{equation}
  b_q(\bm{v})=\mathbb{I}\{\langle\bm{r}_q,\bm{v}\rangle\geq0\},
  \qquad
  d_{\mathrm{H}}(\bm{b}_i,\bm{b}_{i_k})
     =\sum_{q=1}^{m}\mathbb{I}\{b_{i,q}\ne b_{i_k,q}\}.
  \label{eq:rs-eta-hash}
\end{equation}
在各向同性随机超平面和非零向量下，单个位的不同概率与向量夹角成正比，因此汉明距离可近似反映方向差异。选出的$K$条行为进入精细注意力；点击损失更新共享嵌入及预测网络，更新后的嵌入又会影响下一次指纹构造。

假设背包候选指纹为1101，三条历史为1101、1001、0010，则距离为0、1、4。保留两条时，前两条进入精细计算，第三条被排除。若第三条的向量恰好因有限随机投影而被误分，后续注意力无法补回。共享表示缓解了检索与评分各用一套训练依据的问题，但离散符号和Top-$K$不因此成为处处可微运算。指纹更新、长历史扫描与精细计算都需要计入成本。相对SIM，ETA不再依赖类目或单独维护MIPS索引，改用嵌入派生的二进制指纹做按位比较，成本更低；代价是有限位投影带来近似误分，且漏检同样不可由注意力补回。

随机投影矩阵在训练开始时固定，点击梯度更新被选物品的原始嵌入、目标注意、短期分支和预测网络，不更新符号比较或Top-$K$。更新后的物品向量再生成新指纹，因而“端到端”指共享任务训练的表示，不是对汉明距离求连续梯度。服务可把指纹随嵌入表保存，按位异或并统计不同位，再取距离最小的行为；固定机器字内的按位计算便宜，但总扫描仍随历史和候选数增长。

原文在Taobao与一个匿名大型移动应用上评估，工业训练为两周、测试为随后一天，使用最近48条短期和1,024条长期行为。工业表5中ETA加时间的AUC为0.7373、请求推断19毫秒，SIM硬检索加时间为0.7338、21毫秒；消融表7中直接对全1,024条作注意得到0.7371、35毫秒，SimHash选48条得到0.7361、19毫秒，内积选48条得到0.7368、32毫秒。这里存在可测的近似损失，不是哈希保留全部效果。线上相对不含长期分支的DIN基线报告CTR提高6.33\%、GMV提高9.7\%，且已用于实际推荐；原文未公开应用身份，不能自行补成某个已知平台\scite{rs-r76}

\paragraph{SDIM：多哈希的长历史聚合}
\label{sec:rs-sdim}

ETA找到历史后仍逐条计算注意力，候选数量很大时逐条扫描的成本依然突出，可以进一步把候选相关聚合转成桶访问。\term{SDIM}进行多轮随机投影，每轮用若干位组成签名，将与候选同签名的行为向量相加，归一化后跨轮平均，形成长历史兴趣\scite{rs-r77}

\textbf{模型结构。}第一步，行为侧做多轮随机投影，每轮用$b$位组成签名；第二步，把与候选同签名的行为向量相加为桶向量；第三步，归一化后跨轮平均，作为长历史兴趣。设第$r$轮使用$b$位签名$h_r$，相应桶向量为
\begin{equation}
  \bm{B}_{r,c}=\sum_{k:h_r(\bm{v}_{i_k})=c}\bm{v}_{i_k},\qquad
  \bm{h}(i)=\frac{1}{R}\sum_{r=1}^{R}
       \operatorname{norm}_2(\bm{B}_{r,h_r(\bm{v}_i)}).
  \label{eq:rs-sdim-buckets}
\end{equation}
这里$\operatorname{norm}_2(\bm{z})=\bm{z}/\lVert\bm{z}\rVert_2$，空桶按零表示处理。桶内聚合可在行为侧维护，候选只需计算签名并取得相应桶表示。长历史兴趣与短期兴趣和其他特征一起接入CTR网络，用点击交叉熵训练。

若两个非零向量夹角为$\theta$，独立随机超平面下，$b$位全部碰撞的概率为$(1-\theta/\pi)^b$。夹角为$60^\circ$时，两位签名的碰撞概率为$4/9$，四位降为$16/81$。签名更宽减少不同方向混入，也会漏掉更多有一定相似性的事件；增加轮数可改变这种取舍。假设某轮同桶向量为$(1,0)$和$(0.8,0.2)$，其和$(1.8,0.2)$在归一化后约为$(0.994,0.110)$。

同桶不是经核实的相关关系，多轮平均也不是全量softmax注意力的精确重写。它近似的是由碰撞概率诱导的加权方式。与ETA选出$K$条再逐条注意不同，SDIM不再挑选单个事件，而是把同桶向量整体求和、跨轮平均，候选侧成本进一步降到签名查表；代价是桶内细粒度差别被求和掩盖、不同方向可能碰撞混入。行为哈希与在线预测解耦节省了重复工作，但桶的版本、历史窗口及嵌入变化仍需维护。

训练每轮在当前嵌入上重新计算桶归属，固定随机哈希不接受梯度；同桶向量的求和、归一化和预测头仍可把点击梯度传给选中嵌入。原配置用48个哈希方向、每3位组成一个签名，签名宽度与轮数由验证选择。上线把行为侧桶表示计算移到独立服务器，与候选取得并行，CTR服务器读取固定长度桶向量；这减少关键请求路径中的重复工作，但没有消除行为服务器的计算和传输。

Meituan搜索工业实验用14天约100亿实例训练、随后两天评价，长期历史1,024条、短期50条。表3中SDIM的AUC为0.7044，全量长期DIN为0.7049，ETA为0.6906；SDIM相对全量DIN报告11.4倍训练与推断速度，不应将略低的AUC写成全面胜出。宽度消融中3位签名为0.7044，1位为0.6911、10位为0.6923，过宽与过窄都损害效果。14天、每组10\%搜索流量A/B中，2,000条SDIM相对原短期模型CTR提高2.98\%、访问购买率提高2.69\%，请求增加约1毫秒，随后部署主流量。该毫秒主要来自服务器传输，不能用只计算哈希的复杂度代替实际延迟\scite{rs-r77}

\paragraph{TWIN-V2：层次压缩的长历史选择}
\label{sec:rs-twin-v2}

TWIN一致度量仍需访问大量原始事件，历史到十万条后扫描与缓存压力继续增长。\term{TWIN-V2}先把相似历史压成簇，再在簇层面检索和精细聚合，以控制生命周期历史的规模。原任务是短视频CTR：先依据播放完成比例划分行为，再递归聚类，并限制每个簇的最大大小，避免大量相似记录挤在一个过大的簇中\scite{rs-r79}

\textbf{模型结构。}第一步，离线按播放完成比例分组并递归k-means聚类，把相似历史压成簇，限制每簇最大大小；第二步，每个簇形成一个虚拟物品：视频时长、用户播放时长等数值字段取簇内均值，视频ID、作者ID等类别字段取嵌入最接近簇中心的代表物品的字段，而不是保存全簇类别组成；第三步，模型继续使用TWIN式的物品自身特征与交叉特征分解，在线两阶段共享簇级相关性，并根据簇大小调整注意；第四步，把选中的簇表示聚合为兴趣，接入CTR预测，以曝光点击训练。代表字段会丢失其他成员类别差异，后面的数量修正不能将这些差异补回。

压缩必须保留数量的作用。若簇$r$的相关性为$a_r$、大小为$n_r$，两阶段使用调整分数$a'_r=a_r+\log n_r$，精细权重因而正比于$n_r e^{a_r}$。两个相同相关性的簇分别含10个和1个视频时，权重为$10/11$和$1/11$，而不是各占一半。但数量不能恢复具体发生顺序。即使两个簇的平均时长相同，前者可能由一半完整播放、一半快速跳过组成，后者可能全部为中等时长，因而原方案先按播放完成情况分组。

在商品教学类比中，十次相近背包浏览可由一个簇代表，在线访问从十条变为一个簇；这一映射仅解释压缩机制，不把视频的播放完成字段硬套成商品点击字段。与TWIN逐事件选择相比，TWIN-V2把检索对象从原始事件换成约十分之一的簇，扫描量进一步下降；代价是簇代表丢失成员类别差异、具体顺序与异质构成，聚类不可微且需随嵌入更新重做。聚类和物品表示更新会改变簇，簇版本须与在线特征协调。扩大覆盖历史与保留细节形成取舍，不能从“纳入更多记录”推断压缩无损。

离线先按播放完成比例分成5组，在每组递归运行k-means，簇数随待分组规模自适应，直到每簇不超过20条。聚类使用推荐表示，却不是CTR反向传播中的可微层；随后CTR训练更新簇字段嵌入、共享相关性及预测网络。线上保留簇代表、数值均值、簇大小和版本，粗阶段从簇中选100项、细阶段聚合。原部署还保留TWIN近期分支，与生命周期簇分支互补，不是完全以簇取代近期原始行为。

论文主实验最多10万条原始历史，压到约1万簇级输入，对照两阶段方法读取最近1万条事件；因此同时改变了覆盖时间与粒度。连续5天各23小时训练、1小时测试，表3的AUC从TWIN的0.7962升至0.7975，GAUC从0.7336升至0.7360。自适应聚类表4的簇内平均余弦相似度为0.802、每用户聚类0.750秒，平衡二叉方案为0.765、1.010秒；删除任一阶段的簇大小修正均降低效果。线上三个入口相对TWIN观看时长提高0.672\%至0.800\%，多样性指标增益则从0.005\%到0.740\%不等，已部署约4亿日活主流量。设计可扩到百万事件不等于这些离线表已用百万条完成比较；多样性指标定义也不够细，不能直接换称目录覆盖率\scite{rs-r79}

\paragraph{LONGER：合并与查询压缩的长历史建模}
\label{sec:rs-longer}

SIM、ETA、TWIN都在进入主网络前用检索或离线聚类缩小历史，压缩发生在网络之外。\term{LONGER}把一部分压缩放进网络：邻近行为先通过TokenMerge成组合并，可用局部InnerTrans处理组内交互；候选及用户等全局特征形成全局词元。首层用全局词元与少量选定行为作为查询，读取较长的合并历史，后续自注意层继续编码同一短序列，并不逐层再次减少位置数\scite{rs-r80}

\textbf{模型结构。}第一步，相邻$g$个事件通过TokenMerge合并为组表示，组内可用InnerTrans处理局部关系；第二步，候选、用户等全局特征形成全局词元；第三步，首层以全局词元和少量选定行为为查询，读取全部合并历史，查询—历史关系数量由$n^2$降为$mn$；第四步，查询输出进入固定长度的后续自注意与前馈层，不再逐层减位。可以用张量形状理解这种分工：原历史长$n$，相邻$g$个事件组成一组后，组表示的数量约为$n/g$。组内变换处理局部事件关系，跨组层则处理压缩后的表示；首层查询数为$m<n$时，上述计数仅说明注意关系的规模，合并后的宽度、投影和前馈成本仍需另计。

例如，将“登山包、水壶、相机、镜头”分成两个相邻组，局部计算仍可使用每组内两个事件，后续看到的是两个组级表示；若先按候选检索只保留相机和镜头，登山组则根本不会进入精细网络。两种压缩保留的信息不同，不能仅比较最终都是两个输入位置就认为等价。与SIM式先检索再编码相比，LONGER用网络内合并加少查询读取保留全部历史的压缩视图，而不是按候选丢弃尾部事件；代价是合并不可反解、局部顺序只以压缩形式保留。

模型通过候选响应的二元交叉熵训练输入投影、局部处理与后续序列编码，输出点击或转化概率；原论文的主要离线实验是广告CVR，而非只预估点击。合并若采用拼接，本身没有待学习的权重，但后续投影与InnerTrans有。因果注意力限制各位置读取历史的范围，全局词元如何放置也必须满足相应可见性约定。这里的“因果”指计算掩码，和通过干预识别推荐效果的因果推断含义不同。

合并后的组表示一般不能反解每条原始记录，局部顺序也可能只以压缩形式保留。缓存能否复用取决于候选是否已经影响缓存位置、历史是否追加以及参数是否变更。因而长上下文能力要同时检查信息粒度、训练范围与服务预算，不能仅用最大输入长度衡量。

更具体地，原实现加入位置嵌入及事件距请求的时间差，合并还增加向量宽度。首层查询含全局词元和最近$K$个合并位置，键值来自全部输入；查询输出以后进入固定长度的自注意与前馈层。服务掩码阻止历史位置读取当前候选，使历史键值可跨候选复用，逐候选只计算候选全局词元对缓存的读取。混合精度、重计算和稠密稀疏同步更新属于训练执行选择，不改变监督标签；历史追加、内容或模型版本变动仍需刷新相应缓存。

原文以2024年10月16日至2025年2月23日的52亿Douyin广告样本评估CVR，前123天训练、后7天测试，使用48块A100。表1中LONGER的AUC为0.85290、LogLoss为0.47103，Transformer为0.85111、0.47293。表2中不合并配置需$3.73\times10^9$ FLOPs，8词元拼接为$3.03\times10^9$，再加入InnerTrans为$3.52\times10^9$；在后者基础上把250查询改成100个近期查询，降至$1.91\times10^9$，AUC仅略低于250查询配置。随机可学习查询100个却降至0.84946，说明查询来源也影响结果。Douyin电商A/B中直播订单/用户提高7.9222\%、GMV/用户提高6.5404\%，短视频分别提高4.6125\%、5.2771\%，并报告多场景部署。广告ADSS、ADVV另有收益，但原文未给充分指标定义与测试周期，不能把它们擅自解释为收入，也不能由电商结果推出任意场景的同等增益\scite{rs-r80}

\paragraph{STAN：时间与空间条件下的地点匹配}
\label{sec:rs-stan}

前述历史模型只沿时间顺序聚合物品，下一地点还受到空间关系影响，只有顺序仍不够。\term{STAN}以用户、地点和签到时间为输入，分别建立历史事件之间、候选地点与历史之间的时空间隔，用两层注意完成轨迹内更新与候选匹配，避免过早将整个轨迹压成一个向量\scite{rs-r102}

\textbf{模型结构。}第一步，按经纬度计算历史地点两两距离、按事件差值计算时间，形成历史内部的时空矩阵，并计算候选—历史的距离与时间差；间隔通过连续编码或相邻离散表示的插值进入注意关系。第二步，第一层注意把这些历史内时空关系加入事件表示，更新轨迹中各事件；第三步，第二层让候选与全部更新后的事件匹配，把内容兼容性与候选时空关系结合，经注意及可学习聚合得到地点得分。候选—历史时间差在原设定中用目标访问时刻计算。因此，应用若由用户指定“今晚去哪里”，可以使用这个已知时刻；预测一个时间尚未知晓的下一签到时，不能直接读取数据集中未来签到的真实时间。后者需要重新定义可用输入，并重新训练和评价。

用一个明确指定目标时刻的教学例子检查矩阵。为便于手算，以局部平面千米坐标代替经纬度球面距离：三个历史地点为$(0,0),(3,0),(3,4)$，签到时刻为8、9、11时，候选地点为A$(0,4)$、B$(6,4)$，请求目标时刻为12时。两组历史关系为
\[
  \bm{D}_{\mathrm{history}}=
  \begin{pmatrix}0&3&5\\3&0&4\\5&4&0\end{pmatrix},\qquad
  \bm{T}_{\mathrm{history}}=
  \begin{pmatrix}0&1&3\\1&0&2\\3&2&0\end{pmatrix}.
\]
候选—历史距离两行分别为$(4,5,3)$与$(\sqrt{52},5,3)$，时间两行均为$(4,3,1)$小时。即使A、B距最后地点都为3千米，它们与更早轨迹的关系仍有差别。把第二次签到从9时改到10时，会改变历史时间矩阵，并使候选时间行变为$(4,2,1)$。

轨迹聚合与候选匹配逐层消费上述矩阵后输出地点得分。为了只隔离距离输入的作用，假设一个简化评分头固定使用$s_i=-\sum_k d_{i,k}$，则A、B分别为$-12$和约$-15.21$；这是完整网络中的一种教学比较函数，不是STAN保证学习到的距离规则。实际模型以真实下一地点作正例，每步重新采样一定数量的非目标地点，用sigmoid交叉熵训练。服务时去掉训练采样器，对适用地点目录计算分数。

签到并不覆盖全部实际活动，GPS还可能有噪声；距离接近也不代表道路可达或正在营业。模型输出的是访问倾向，候选资格需另查。时间与空间条件改变输入和评价任务，不能把得分提高解释为自动满足真实出行要求。与只做序列匹配的历史模型相比，STAN额外显式建模历史内部与候选—历史两层时空关系，而不是把位置当作普通物品标识。

训练前缀依次展开，每步采样非目标地点，使负例随训练变化；默认每个正目标取10个负地点，Adam同时更新用户、地点、时间、间隔表示和两层注意、匹配聚合参数。候选匹配输出是地点分数，并不保证不同地点分数和为一；服务取消训练负采样器，枚举或访问适用地点。原训练的目标时刻可见条件与本文已知请求时刻的例子必须对应，不能只因前缀掩码合法就忽略候选时间输入。

Gowalla、东京、纽约和新加坡签到实验以每人末次签到测试、倒数第二次验证，最大轨迹100条。表2纽约Recall@10为0.5962，GeoSAN为0.5267、TiSASRec为0.5020；表3去掉时间间隔为0.5814、去掉空间间隔为0.5860，两者与平衡采样全部去掉为0.5094。报告因此支持时空关系、两层匹配与采样共同作用，不能把全部差异归给空间距离。论文也说明部分基线按各自轨迹切片，并取消ARNN的语义知识，未严格固定所有信息与窗口。其实际应用对象是地点推荐日志，不是已经验证营业、道路与目标时刻未知的在线出行规划\scite{rs-r102}

这些候选相关模型还需与信息保存的代价一起比较。表~\ref{tab:rs-history-tradeoffs}固定“历史已经取得、候选需要评分”这个任务，列出请求时访问什么、行为到来时维护什么。它不是原论文性能排名：原始历史范围、标签、候选数和硬件没有统一，不能以各自报告的毫秒直接排序。

\begin{table*}[htbp]
  \centering
  \small
  \begin{tabularx}{\linewidth}{@{}lXXX@{}}
    \toprule
    同行方案 & 保存及更新对象 & 请求侧主要计算 & 适用条件与主要损失 \\
    \midrule
    DIN、DIEN、DSIN & 原始事件或会话编码 & 候选激活、演化或会话聚合 & 历史可逐候选读取；长窗重算昂贵 \\
    MIMN & 固定槽记忆及归纳状态 & 读记忆、候选预测 & 事件更新及时；难恢复旧细节 \\
    SIM、ETA、TWIN & 原始历史及类别索引、指纹或投影 & 选事件后精细注意 & 稀疏相关证据可筛选；漏检不可补回 \\
    SDIM & 多轮历史桶表示 & 候选签名、桶读取与聚合 & 候选很多且容许近似；碰撞混入与漏桶 \\
    TWIN-V2 & 簇代表、均值、大小与近期历史 & 簇选择及聚合 & 生命周期重复较多；丢失成员类别与顺序 \\
    LONGER & 合并输入及合法历史键值 & 首层少查询、固定短序列编码 & GPU与缓存可用；合并及查询压缩损失 \\
    PinnerFormer、TransAct & 批量长期向量及即时事件 & 共享长期特征、逐候选近期编码 & 多日偏好与当前变化并存；维护两条时间线 \\
    \bottomrule
  \end{tabularx}
  \caption{候选响应任务中的历史保存与访问取舍。参数、事件、索引、桶簇和键值缓存是不同状态，更新其中一项不自动更新其余项。}
  \label{tab:rs-history-tradeoffs}
\end{table*}

若请求包含几百个候选，候选无关历史编码和内容向量的复用尤其重要；若关键线索只在很少旧事件中，SIM、ETA、TWIN的细粒度选集更便于追查，但要测漏检。生命周期有大量重复视频时，簇和桶可减少访问，不过簇代表不能用于逐条事实证明。需要新收藏立即影响判断时，实时分支或增量状态能补批量向量的时效缺口；它们也要求去重、事件可见时间和模型版本一致。选择机制应先确定允许丢掉哪些信息，再比较相同预算下的任务质量。

% Outline: 4.3.1.3
\subsubsection{基于语义内容的评分}
\label{sec:rs-semantic-scoring}

物品名称、说明和用户表达可以补充ID之外的含义。语义向量可作为字段送入上述评分器，也可让语言模型直接输出评分、接受判断或候选次序。后一种接口要特别检查身份、提示和解析协议：输出了一段合理文字，不等于已经比较了整个目录，更不等于概率经过校准。

% Outline: 4.3.1.3.1
\paragraph{P5：文本化响应与推荐判定}
\label{sec:rs-p5-scoring}

纯ID嵌入的评分器只能区分“见过哪个物品”，读不到物品名称、说明和用户自然语言表达；而把评分、序列推荐、解释、评论等任务各自训练一个回归头，又会让同一套语言知识无法复用。\term{P5}（Recommendation as Language Processing，统一预训练、个性化提示与预测范式）把不同推荐任务都表达为文本输入和文本目标，用同一个序列到序列语言模型学习，使评分数字、肯定否定判断、下一个物品编号都成为词元预测\scite{rs-r21}

\textbf{模型结构。}第一步是任务模板化：当前候选评分时，输入文本由任务提示、用户标识、用户已知偏好、物品标识及可用描述拼接而成，目标文本可以是评分数字（如“4”），也可以是“是/否”判断。第二步是词元级建模：模型按自回归方式预测目标文本，训练目标是目标词元的负对数似然，而不是对输出数值施加同一个回归均方误差；评分数字被当作离散词元序列处理，模型学到的是“在该提示下该词元出现的概率”。第三步是服务解析：限定合法输出集合，把词元解码为数值或判断后再比较候选；若输出“很好”或超出量表的“8”，应进入明确的无效输出处理，而不能静默当成高分。不同候选必须遵循相同模板和解码协议。

教学模板为“用户$u$的已知偏好为……；预测其对物品$i$的1至5分评分”，监督文本为“4”。这里的关键区别在于：词元“4”的概率受提示措辞、词表排序和相邻词元影响，它不是经过校准的响应概率；生成一个“是”只表示模型作出了一次文本判定，不能直接换算成CTR。

P5还包括序列推荐、直接推荐、解释和评论任务，这些任务的输出语义不同（下一个物品编号、理由文字、评论续写），不能共用一个数值尺度。训练按各任务文本的词元似然联合进行，服务时按任务选择模板与合法词元集合。用户ID、物品ID及模板切分还可能把未来行为或目标信息泄漏进提示，须在第~\ref{chap:rs-evaluation}~章的协议下检验。

与下面TALLRec相比，P5从预训练语言模型出发用提示覆盖多任务，TALLRec则聚焦给定目标的肯定/否定判定并做参数高效适应；两者都把语言流畅与排序有效性分开评价：输出一段通顺的理由，不等于模型比较了整个目录，更不等于词元概率经过校准。

% Outline: 4.3.1.3.2
\paragraph{TALLRec：语言模型推荐判定}
\label{sec:rs-tallrec}

通用语言模型虽读过大量商品文本，却未必把历史偏好与指定候选对应起来：直接询问“用户是否喜欢这个物品”，模型可能凭物品自身热度作答，而忽略用户历史里的具体证据。\term{TALLRec}把用户历史中的喜欢、不喜欢等证据及目标物品写入推荐指令，用明确的肯定、否定反馈对模型做参数高效适应，使判断落到给定目标上\scite{rs-r22}参数高效适应只改变训练哪些参数（如低秩支路或提示参数），冻结主体语言模型，当前任务仍是“对给定目标作判断”，不是让模型自行访问目录。

\textbf{模型结构。}第一步构造指令：固定一段偏好历史（喜欢/不喜欢的物品及理由），把目标物品名称填入模板，形成“是否推荐该物品”的判定输入。第二步监督判定：训练为每个目标提供真实任务标签，模型在约定的肯定、否定词元上学习对应概率。第三步恢复比较分数：服务时读取允许的回答，或在肯定、否定词元之间恢复比较分数。若以两个词元概率归一化为
$q=p_{\mathrm{yes}}/(p_{\mathrm{yes}}+p_{\mathrm{no}})$，
应明确这是所选解析协议，还要在目标数据上校准和评价，不能把任意自然语言回答直接换成CTR。

教学输入固定同一段偏好历史，分别询问通勤背包与摄影镜头：历史证据不同，两个目标的肯定词元概率应给出不同判断，而不是对所有目标输出同一套热门物品偏好。训练只更新被开放的参数，主体模型保持冻结；服务时逐候选构造指令并解码，调用次数随候选数增长。

候选标题有时不足以唯一确定物品，内容也可能掩盖用户过去行为的具体对象。应检查身份映射、提示长度、标签数量及领域变化，分别报告低数据量实验的范围。自然语言理由可辅助阅读，但不能证明模型识别了真实偏好原因；与P5相比，TALLRec不追求覆盖序列推荐、解释等多任务，而是把判定接口做窄、做稳。

% Outline: 4.3.1.3.3
\paragraph{RankGPT：指令候选重排}
\label{sec:rs-rankgpt}

P5与TALLRec逐候选作点式判定，模型每次只看到一个目标，无法直接比较候选之间的相对优劣，且调用次数随候选数线性增长。\term{RankGPT}让语言模型直接比较一组文本：把需求和已有候选编号交给模型，要求返回候选排列；长列表则通过重叠滑动窗口逐段重排\scite{rs-s14}原论文主任务是搜索相关性重排，个性化推荐迁移还需提供用户偏好证据，不能把搜索实验当成推荐效果。

\textbf{模型结构。}第一步是列表构造：把用户需求（推荐场景下为偏好历史与当前意图）与候选编号文本一并写入提示，候选以$[1],[2],\ldots$编号。第二步是窗口内排列：模型在每个窗口内返回候选次序；长列表采用由尾部向前处理的重叠窗口，允许原来靠后的候选经多轮比较逐步向前移动，窗口宽度和重叠步长共同决定调用次数与可比较范围。第三步是全局恢复：把各窗口局部排列合并回原候选池位置。

假设候选编号为$[1],[2],[3],[4]$，输出$[3]>[1]>[4]>[2]$只表示当前比较的次序，不附带分数。实现必须检查编号是否重复、缺失或越界，规定错误输出的回退，并保留每次窗口输入与恢复后的全局位置。

在同一候选池上改变输入顺序，是检查位置敏感性的一个直接办法：若次序随编号排列翻转，说明模型受提示位置影响而非稳定比较。还应把候选覆盖与重排质量分别评价：模型不能把未提供的好物品排到前面，重排只在召回给定的池内生效。若把重排结果蒸馏给较小评分器，需明确教师次序、训练候选及蒸馏损失；阶段预算在第~\ref{sec:rs-stage-connection}~节讨论。

与09章搜索侧的RankGPT相比，机制（列表提示、重叠窗口、排列解析）一致，差别在输入证据：搜索侧比较查询与文档相关性，推荐侧必须额外给出用户偏好历史，否则模型只能按物品通用热度排序。语言流畅不等于排序有效，窗口合并规则、位置偏置与错误解析都会改变最终次序。

% Outline: 4.3.2
\subsection{评分模块的组合与统一}
\label{sec:rs-scoring-organization}

字段和历史都能影响分数，但还要决定两者何时相遇。先把历史压成一个向量再做字段交互，计算简单，却可能提前丢失候选需要的细节；逐层交换信息或放入共同骨干，可以保留更多交互机会，也改变复用条件。下面比较完整评分器的组织，不是要求服务依次执行这些模型。

% Outline: 4.3.2.1
\subsubsection{基于分离模块的融合评分}
\label{sec:rs-separate-fusion}

一种直接方案是用DIN得到候选兴趣，拼接用户、候选和情境字段，再交给DCN-V2；也可以让序列编码器先产生用户表示，作为RankMixer的一个输入单元。两者保留历史与字段模块各自的计算结构，但通过共同响应损失联合训练。若历史表示依赖候选，在线仍需按候选重新计算相应部分。

分离模块不必只在末尾拼接。FinalMLP融合两条字段分支，InterFormer让字段和历史模块逐层交换摘要。比较时应明确融合对象、传递信息和发生位置，并固定历史长度、候选数及参数预算；否则更长历史带来的收益容易被归到“融合结构”上。

% Outline: 4.3.2.1.1
\Needspace{90mm}
\paragraph{FinalMLP：互补分支的融合评分}
\label{sec:rs-finalmlp}

双流CTR模型通常让两条分支各自加工后只在最终logit相加，分支之间在网络内部从不交换信息，两条同构MLP因此容易学到高度重复的表示。\term{FinalMLP}用上下文门控分别调整两分支的输入，再以多头双线性聚合两个分支的输出，使两条分支在输入端和输出端都形成互补\scite{rs-r53}

\textbf{模型结构。}第一步是门控选输入：设字段拼接为$\bm z$，两个门控为$\bm g_1,\bm g_2$，门控由上下文产生并逐维缩放，两分支分别读取被筛选后的输入，$\bm h_j=\operatorname{MLP}_j(\bm g_j\odot\bm z)$。第二步是双线性融合：融合同时包含两侧线性贡献与双线性作用，单头可以写成
$s=b+\bm w_1^\top\bm h_1+\bm w_2^\top\bm h_2+\bm h_1^\top\bm W\bm h_2$。
第三步是多头分块：对表示分块后分别做双线性，减少完整双线性矩阵$\bm W$的参数量，各头结果再汇总。

若两个教学分支输出为$(1,0.5)$和$(0.2,1)$，$\bm W$取单位矩阵，则双线性贡献为0.7；假设偏置与线性项合计为$-1$，最终概率为$\sigma(-0.3)\approx0.426$。点击标签通过同一交叉熵更新门控、两分支及融合层；门控随上下文变化，服务时每个候选都要重算门控与双线性，不能把某一分支输出预存后跳过融合。

FiBiNET交互的是字段向量，FinalMLP此处交互的是两个已加工分支的输出，不能只因都用了双线性项便视为同一结构；与InterFormer相比，FinalMLP在两条平行分支的末端融合，不做逐层信息交换。

% Outline: 4.3.2.1.2
\paragraph{InterFormer：字段与历史的逐层融合评分}
\label{sec:rs-interformer}

只让字段指导历史聚合（DIN式候选注入历史注意），没有利用历史反过来修正字段计算：第一层算定的字段表示不会再被发现的历史偏好更新。\term{InterFormer}保留字段交互模块与序列模块两套结构，让两者在各层通过桥接模块交换摘要，形成双向逐层融合\scite{rs-r56}

\textbf{模型结构。}第一步，字段交互模块产出字段摘要，送入序列模块的注意与个性化前馈，使历史读取随当前字段调整。第二步，序列模块把全局、多方面及近期摘要回送字段交互模块，使字段计算吸收历史证据。第三步，逐层重复上述交换：字段表示不再是第一层的固定结果，而在后续层结合回传的历史摘要重算。逐条历史和字段表示仍分别保留，不必把两者全部压成同一种单元；最终字段表示接响应头，以曝光点击交叉熵共同训练两侧模块和桥接。

例如当前字段指出候选为通勤背包，第一层字段摘要使序列模块关注通勤事件；序列侧又发现近期多次查看大容量商品，将相应摘要交回字段模块。第二层字段交互便可结合容量、价格与当前兴趣重新计算，而不是一直使用第一层的固定历史摘要。

桥接摘要仍是压缩，粒度和宽度会限制可传递的信息；增加层间交换也会增加计算依赖，序列侧计算若依赖候选，仍须逐候选重算。它与下文OneTrans的区别在于保留两套模块并交错融合，而不是共用一套逐层骨干；两者都只解决候选评分内部的组织问题。

% Outline: 4.3.2.2
\subsubsection{基于联合骨干的评分}
\label{sec:rs-joint-backbone}

另一种组织方式是让字段和历史持续进入共同的逐层计算。联合骨干并不要求所有参数共享，也不要求任意输入彼此可见。MTGR保留业务交叉并按用户聚合，OneTrans将行为及字段组织为共同序列，MixFormer让字段查询反复交互并读取历史；它们都仍需要外部候选。

\Needspace{110mm}
% Outline: 4.3.2.2.1
\paragraph{MTGR：交叉特征与行为联合评分}
\label{sec:rs-mtgr}

把输入改成行为序列时，若丢弃已有的用户—候选交叉统计（历史共现、即时响应等），序列编码器可能同时丢掉长期有效证据；而纯CTR交叉模型又不便于按用户聚合、复用窗口前的计算。\term{MTGR}保留用户字段、历史、即时行为及候选交叉特征，将它们映射到共同维度，以HSTU式编码器产生各候选响应\scite{rs-r40}MTGR的作者释义是Meituan Generative Recommendation，这里的输出是给定候选的分数，并非直接解码整张推荐列表。

\textbf{模型结构。}第一步，把用户字段、静态历史、窗口内即时行为与候选交叉特征分别映射到共同维度，不同语义组采用相应归一化。第二步，按用户聚合同一窗口的候选，窗口之前的静态历史和用户计算只算一次、供窗内各候选复用。第三步，动态掩码控制可见性：窗口前的静态字段、历史可以被候选读取；窗口内的新行为只对时间更晚的候选可见；候选之间隔离，不能借另一条样本泄漏当前不可用的信息。HSTU基础计算见第~\ref{sec:rs-hstu}~节。

例如上午10点候选$a$、下午2点候选$b,c$被放进同一训练窗口，用户中午12点点击相机。该点击可以供$b,c$评分，不能供$a$评分；$b$之后才得到的响应也不能作为$c$的输入，除非$c$的真实请求时点已经观测到它。各候选仍按自己的曝光标签计算损失。用户计算的复用减少重复工作，但候选表示与输出不会因聚合而完全不花成本；应按整次请求计量。

与OneTrans把行为与字段堆成一条因果序列不同，MTGR保留交叉特征并按用户窗口组织可见性；它仍以给定候选为前提输出分数，不是直接生成整张推荐列表。

% Outline: 4.3.2.2.2
\paragraph{OneTrans：特征交互与行为序列联合评分}
\label{sec:rs-onetrans}

历史先汇总成固定向量，会限制后续字段能选择的行为细节；而字段交互模块与序列模块分栈训练时，两者只在末端拼接，多层之间也不交换信息。\term{OneTrans}把行为单元放在字段单元之前，由同一Transformer堆栈处理，让历史与字段在同一逐层计算中交互；行为侧共享查询、键、值和前馈参数，不同字段单元则保留各自参数\scite{rs-r41}通用注意计算见\BookChapterRef{chap:transformer}{模型卷的“Transformer”章}，本节关心混合参数和可见范围。

\textbf{模型结构。}第一步是单元排布：四条行为$b_1,b_2,b_3,b_4$在前，随后是用户字段$u$和候选字段$i$，全部进入同一堆栈。第二步是因果可见：掩码使$i$读取全部此前行为及$u$，使$u$读取此前行为，但行为不能反向读取后来的候选。第三步是参数分工：行为单元共享注意与前馈参数，字段单元保留专属参数；候选输出已经含有历史和字段的多层交互，可接CTR、CVR等响应头，以各任务对应标签联合训练。多行为输入若按意图而非时间排列，掩码含义也应随其组织方式说明；不能将任意排列都叫时间因果关系。

金字塔计算逐层减少发出查询并保留输出的行为单元。例如某层只保留$b_3,b_4,u,i$的查询，但该层键和值仍来自$b_1$至$i$的完整当前输入；下一层才接收缩短后的输出。这与一开始删除$b_1,b_2$不同：较晚单元已经有机会读取它们。若当前层输入长度为$L$、保留查询数为$L'$，注意乘加约为$O(LL'd)$，前馈只处理保留输出。

历史不能读取候选，使同一请求中历史键值可在多个候选间复用。跨请求复用还依赖追加式历史、相同排序方式及模型版本；事件修改、候选相关子序列或改变位置编码，都可能破坏直接复用条件。新增行为的键值投影可以只做增量部分，但新查询读取旧缓存仍有代价。OneTrans统一评分内部计算，跨阶段联合学习另见第~\ref{sec:rs-onetrans-v2}~节。

% Outline: 4.3.2.2.3
\Needspace{95mm}
\paragraph{MixFormer：特征交互与历史聚合联合评分}
\label{sec:rs-mixformer}

联合评分不一定把所有历史和字段放进同一自注意序列：那样历史与字段互为键值，顺序与成本都被拉长，也难以单独复用用户侧计算。\term{MixFormer}将非序列字段组成多个查询头，历史单独作为键值，每层依次进行查询混合、历史读取及输出融合\scite{rs-r42}

\textbf{模型结构。}第一步是查询混合（Query Mixer）：通过头间信息重排及各头前馈交换字段信息，让用户、价格、容量等字段在查询头之间组合。第二步是历史读取：这些查询对历史键值作交叉注意，按当前字段条件选择历史证据。第三步是输出融合：各头前馈把当前字段与读到的历史结果融合，堆叠后的表示交给多任务响应头。

沿背包候选看两层计算：第一层查询已经包含用户、价格、容量等字段，读取到通勤相关历史；第二层的查询混合继续组合这些结果，新的查询可以强调与容量或近期购买相关的事件。与固定一次历史汇总相比，后层读取条件会变化；与OneTrans相比，历史主要作为被查询的序列，交互位置和参数组织不同。

用户—候选解耦变体进一步限制信息流，使候选侧可以读取用户侧，而用户侧不接收候选信息，从而复用同次请求的用户计算。若允许候选先修改用户查询，就不能再把修改后的结果直接供其他候选使用。响应标签、头数、历史长度和稠密容量需共同记录；联合扩展的收益应连同论文的数据和预算检验，不能仅凭名称推断网络结构或全流程整合范围。

图~\ref{fig:rs-scoring-organization}强调三个不同的融合位置；表~\ref{tab:rs-scoring-reuse}进一步比较可复用部分。共同骨干与共享缓存是相关但不同的选择，后者由信息依赖决定。
\begin{figure*}[htbp]
 \centering
 % 字段与历史的三种融合位置；全版心，约162mm。
\begin{tikzpicture}[
 font=\figurefont\footnotesize,
 box/.style={draw=BookRule,fill=BookPaper,align=center,
   text width=32mm,minimum height=11mm,inner sep=3pt},
 flow/.style={-{Stealth[length=2mm]},draw=BookMuted,line width=.9pt}
]
 \node[font=\figurefont\footnotesize\bfseries,anchor=west] at (-1.8,1.35)
   {末端融合：先独立压缩，再比较候选};
 \node[box,fill=FigureInput!10] (h1) at (0,.55) {原始历史};
 \node[box,fill=FigureInput!10] (f1) at (0,-.85) {用户、候选与情境};
 \node[box,fill=FigureModel!10] (he1) at (4.8,.55) {历史表示};
 \node[box,fill=FigureModel!10] (fi1) at (4.8,-.85) {字段表示};
 \node[box,fill=FigureOutput!10] (o1) at (12,0) {融合及响应头};
 \draw[flow] (h1)--(he1);
 \draw[flow] (f1)--(fi1);
 \draw[flow] (he1.east)--(8,.55)--(o1.west);
 \draw[flow] (fi1.east)--(8,-.85)--(o1.west);
 \node[font=\figurefont\footnotesize\bfseries,anchor=west] at (-1.8,-2)
   {逐层桥接：两套表示通过摘要交换信息};
 \node[box,fill=FigureModel!10] (h2) at (0,-2.85) {历史模块};
 \node[box,fill=FigureModel!10] (f2) at (5.5,-2.85) {字段模块};
 \node[box,fill=FigureOutput!10] (o2) at (12,-2.85) {响应头};
 \draw[flow] (h2.east) to[bend left=20] node[above]{历史摘要} (f2.west);
 \draw[flow] (f2.west) to[bend left=20] node[below]{字段摘要} (h2.east);
 \draw[flow] (f2)--(o2);
 \node[font=\figurefont\footnotesize\bfseries,anchor=west] at (-1.8,-4.3)
   {联合骨干：每层同时安排字段交互与历史读取};
 \node[box,fill=FigureInput!10] (i3) at (0,-5.2) {历史与字段单元};
 \node[box,fill=FigureModel!10] (b3) at (5.5,-5.2) {联合层堆栈\\掩码控制信息可见};
 \node[box,fill=FigureOutput!10] (o3) at (12,-5.2) {候选响应头};
 \draw[flow] (i3)--(b3);
 \draw[flow] (b3)--(o3);
\end{tikzpicture}

 \caption{候选评分中字段与历史相遇的位置。箭头表示信息依赖，图中“联合层”并不表示全部输入双向可见；具体掩码及参数共享仍由各模型规定。}
 \label{fig:rs-scoring-organization}
\end{figure*}

\begin{table*}[htbp]
 \centering\small
 \begin{tabularx}{\linewidth}{@{}lXXXX@{}}
 \toprule
 组织 & 原始输入与交互位置 & 可见范围与参数 & 计算复用 & 输出及主要代价\\
 \midrule
 分离模块融合 & 字段与历史分开加工；末端或逐层摘要融合 & 由模块接口限定；参数分别组织 & 独立用户编码可复用；候选相关聚合须重算 & 响应头；两侧模块及桥接\\
 MTGR & 字段、静态历史、即时行为和候选交叉共同编码 & 动态时间掩码；候选隔离；语义组归一化 & 用户窗口聚合；仍需候选计算 & 多候选响应；窗口大小与掩码\\
 OneTrans & 行为在前、字段在后共同堆叠 & 因果可见；行为共享、字段专属参数 & 请求内历史缓存；跨请求需追加与版本一致 & 响应头；保留查询数与序列长度\\
 MixFormer & 字段查询逐层混合并读取历史 & 查询头参数化；解耦变体限制候选回流 & 解耦后复用请求内用户侧；跨请求另验 & 多任务头；查询数、历史与稠密容量\\
 \bottomrule
 \end{tabularx}
 \caption{评分模块组织与复用条件。所有行均以给定候选为前提；固定信息时点、候选数、标签和预算，才能区分架构变化与更多信息、更多计算带来的收益。}
 \label{tab:rs-scoring-reuse}
\end{table*}

% Outline: 4.3.3
\subsection{观测数据下的评分学习}
\label{sec:rs-observed-learning}

评分器只见到日志中的样本，而日志并非从所有用户—物品对均匀抽取。旧系统决定展示什么，位置影响是否查看，用户点击还可能受流行度信息影响。缺失、误点与恶意反馈又会改变标签含义。处理这些条件时，应先定义事件和可观测信息，再决定重加权、建模或数据治理怎样进入训练。

% Outline: 4.3.3.1
\subsubsection{反馈标签与质量控制}
\label{sec:rs-feedback-quality}

人工相关性判断、用户显式评分、点击、消费和生成样本提供不同监督。人工标注可能只看内容而不知道个体需求；点击要求对象先被展示且被注意到；消费依赖后续追踪；语言模型生成的偏好则是另一个模型的判断。不能把它们统一改成0、1后就认为目标相同。

一份可用样本应能说明请求、候选、展示位置、结果事件及观察截止时点。曝光不等于查看，未点击可能因为没有注意到，点击也可能是误点。若观察窗口为7天，发生1小时的样本仍未成熟；此时“尚未记录购买”与“7天没有购买”不同。延迟回传和广告标签修正见第~\ref{chap:rs-advertising}~章。

质量检查应追踪字段的实际可用时间。最终成交价、事后评分和未来聚合统计不能作为请求时输入；以生成文本补充内容时，也需排除目标标签及未来评论。错误追踪造成的缺失应先修正记录或标记可用范围，不能期待网络自动辨认每种缺失原因。

% Outline: 4.3.3.2
\subsubsection{负例构造与训练采样}
\label{sec:rs-ranking-negatives}

未交互、曝光未点击、人工判为不合适和从目录采样，分别对应不同的负例依据。曝光未点击至少说明进入过展示，但不保证被查看；未交互物品可能从未有机会出现；人工负例的有效范围由标注任务决定。BPR需要偏好比较假设，CTR则需要响应事件及窗口，不能共享一个未经说明的“负标签”。

采样还会改变概率。假设所有正例保留，负例以常数概率$a$独立保留，且除此之外不依赖输入；若真实概率为$p$，采样后概率$q$满足
\[
 q=\frac{p}{p+a(1-p)},\qquad
 p=\frac{aq}{1-q+aq}.
\]
教学例中$a=0.1$、模型在采样数据上预测$q=0.5$，恢复的$p$约为0.0909。输入相关采样时需使用对应的$a(\bm x,i)$；漏掉该修正会把训练分布概率误当成线上概率。

难负例可加强局部竞争，也可能包含真实喜欢但未观察的物品。召回训练面对目录竞争，排序面对上游筛选后的候选；第~\ref{sec:rs-retrieval-training}~节的采样原则可复用，但需要重新说明比较分布。负采样处理的是训练计算与权重，不能替代曝光或位置机制的分析。

% Outline: 4.3.3.3
\subsubsection{曝光与点击选择偏差下的评分学习}
\label{sec:rs-selection-bias}

旧系统只展示过一款背包时，另一款的缺失反馈不能直接当成拒绝。Recommendations as Treatments将反馈是否被观察记为$O$，用观测倾向$o_{u,i}=\mathbb P(O_{u,i}=1\mid\text{已说明的信息})$校正风险\scite{rs-e01}若目标是全体用户—物品对上的平均损失，一种逆倾向估计为
\begin{equation}
 \widehat R_{\mathrm{obs}}
 =\frac1{|U||I|}
   \sum_{(u,i):\,O_{u,i}=1}
    \frac{\ell(y_{u,i},\widehat y_{u,i})}{o_{u,i}}.
 \label{eq:rs-observed-risk}
\end{equation}
观测概率为0.5与0.1的两个样本，其损失权重为2与10；低概率样本代表更多未观察机会，也带来更大方差。

这个估计需要正确的观测机制、足够覆盖及相应可忽略性条件。若某类对象从不被观察，就无法单靠这份日志识别其平均损失；模型外推应与日志支持的范围分开。倾向未知时还需交代估计来源，不能以响应模型自己的点击预测代替观测概率。

广告CVR只在已点击样本中直接取得转化监督，属于进入训练的选择条件，后文用ESMM和ESCM$^2$讨论。观测倾向、位置检查概率、点击概率，以及策略评价中的日志行动概率具有不同事件与分母，必须分别命名。

% Outline: 4.3.3.4
\subsubsection{位置偏差下的评分学习}
\label{sec:rs-position-bias}

即使对象已经展示，低位结果也可能因为较少被查看而缺少点击。令$E$表示查看，$R$表示在当前请求下具有吸引力或相关性，$C$表示点击，$k$为位置。在简化的位置点击模型中，假设点击需要查看及正响应，查看只依赖位置，相关性不再受位置影响，得到
\begin{equation}
 \mathbb P(C=1\mid\bm x,i,k)=e_k r(\bm x,i).
 \label{eq:rs-pbm}
\end{equation}
这个乘积来自生成假设，不是点击的定义。若用户因信任首位而点击原本不喜欢的对象，或者前一结果改变了后续查看行为，简单分解就可能不成立。

图~\ref{fig:rs-position-paths}还表明，旧排序同时影响位置分配，相关性及其代理特征也参与旧排序。因此首位点击率更高，既可能来自更常查看，也可能来自候选本身更相关；直接比较位置均值不能分离二者。随机换位、可比较的历史排序器变更或额外结构假设，才为分解提供证据。
\begin{figure}[htbp]
 \centering
 % 位置、查看与信任的简化关系；正文宽。
\begin{tikzpicture}[
 font=\figurefont\footnotesize,
 box/.style={draw=BookRule,fill=BookPaper,align=center,
   text width=20mm,minimum height=10mm,inner sep=3pt},
 flow/.style={-{Stealth[length=1.8mm]},draw=BookMuted,line width=.9pt}
]
 \node[box,fill=FigureInput!10] (r) at (0,2) {相关性\\及其代理特征};
 \node[box,fill=FigureModel!10] (old) at (3.4,2) {旧排序};
 \node[box,fill=FigureInput!10] (pos) at (6.8,2) {展示位置};
 \node[box,fill=FigureModel!10] (e) at (6.8,0) {查看};
 \node[box,fill=FigureModel!10] (trust) at (3.4,0) {信任};
 \node[box,fill=FigureOutput!10] (c) at (0,0) {点击};
 \draw[flow] (r)--(old);
 \draw[flow] (old)--(pos);
 \draw[flow] (pos)--(e);
 \draw[flow] (r)--(c);
 \draw[flow] (e.south)--(6.8,-1)--(0,-1)--(c.south);
 \draw[flow,dashed] (pos.south west)--(trust.north east);
 \draw[flow,dashed] (trust)--(c);
\end{tikzpicture}

 \caption{位置点击偏差的简化因果关系。实线表示模型中保留的作用，虚线表示可能存在的信任路径；相关性及其可用代理影响旧排序，也影响响应。只观察点击和位置通常不足以唯一分离这些作用。}
 \label{fig:rs-position-paths}
\end{figure}

% Outline: 4.3.3.4.1
\paragraph{Propensity SVM-Rank：基于倾向加权的排序学习}
\label{sec:rs-propensity-svm}

普通Ranking SVM把每次点击当作同等正反馈，低位对象本来更少被查看，其点击在训练中被系统性低估，学到的分数会偏向占住高位的对象。已知检查倾向后，可以让一次低位点击代表更多本来可能被查看的正反馈。\term{Propensity SVM-Rank}对点击对象相对其他候选的排序损失作逆检查倾向加权，再训练带正则的Ranking SVM\scite{rs-s09}若线性分数为$s_i=\bm w^\top\bm\phi(\bm x,i)$，可用成对hinge项$\max(0,1-s_i+s_j)$上界点击对象的名次损失，再乘以$C_i/e_{k_i}$：权重来自检查概率的倒数，位置越低、$e_k$越小，一次点击承担的训练权重越大。

假设同一对象在两个位置的检查概率为0.5与0.1，一次点击的权重分别为2与10；若某个比较对的hinge损失为0.4，贡献为0.8与4。较少被查看的对象因此不会仅因低位而在训练中被系统性忽略。这里把点击对象与其他对象比较，是构造相关性加权名次风险，不是断言所有未点击对象都是真实负例。线上只按学到的特征分数排序。

倾向从哪里来是无偏性的前提。原研究通过少量随机换位标定相对位置倾向；另一项工作从多个历史排序器中寻找同一查询—文档出现在不同位置的记录，提取位置比较证据\scite{rs-s33}后一协议仍依赖位置模型、覆盖和记录可比性，多排序器日志本身不等于随机处理。逆倾向把低概率样本放大$1/e_k$倍，小倾向会显著增加方差；截断权重也会改变无偏性，二者都需报告。与下面Regression-based EM相比，本方法假定倾向已知或已由外部标定，不反过来从点击中估计倾向。

% Outline: 4.3.3.4.2
\paragraph{Regression-based EM：基于潜变量估计的倾向学习}
\label{sec:rs-regression-em}

Propensity SVM-Rank要求倾向已知，而个人搜索中同一查询—文档可能很少重复，既无法逐对统计相关性，也没有随机换位实验可直接标定每个位置的检查概率。\term{Regression-based EM}用特征回归$r(\bm x,i)$在稀疏重复下共享信息，并把查看与相关性当作点击生成中的潜变量，交替估计其后验和更新参数\scite{rs-s32}

在式~\eqref{eq:rs-pbm}的假设下，点击为1时必有$E=R=1$；点击为0时有多种解释（没被查看、查看了但不相关，或两者兼有）。固定当前估计$e,r$，未点击样本的后验为
\begin{equation}
 \mathbb P(E=1\mid C=0)=\frac{e(1-r)}{1-er},\qquad
 \mathbb P(R=1\mid C=0)=\frac{r(1-e)}{1-er}.
 \label{eq:rs-rem-posterior}
\end{equation}
取教学$e=0.5,r=0.4$，两个后验为0.375与0.25。未点击并没有强迫相关性后验变成0。

E步据当前参数计算这些后验，M步以各位置的查看后验更新检查参数，以相关性后验更新回归。原实现从后验采样相关性标签，并用GBDT学习特征回归；所得位置倾向再供后续加权排序器使用。即使两篇个人文档不重复，它们的特征也可通过回归共享统计依据。整个EM离线迭代，线上使用最终排序器，不要求实时执行整套EM。

EM提高所设模型的拟合程度，却不能单独证明查看与相关性的分解唯一。模型表达错误、位置覆盖不足及强信任偏差仍会使倾向失准；它与Propensity SVM-Rank的分工在于后者直接使用外部标定倾向，本方法从点击与特征中估计倾向。

% Outline: 4.3.3.4.3
\paragraph{DLA：基于倾向与相关性互估的排序学习}
\label{sec:rs-dla}

前述方法先估倾向再喂给排序模型，两步误差会沿链条传播：倾向估低会过度放大某次点击，相关性模型的偏差又会反过来污染倾向估计。相关性有助于估计查看机会，查看机会又有助于学习相关性。\term{双重学习算法}（dual learning algorithm, DLA）将两者组成联合训练问题\scite{rs-s10}排序模型读取查询—文档特征，位置模型读取位置，分别在列表内形成归一化分数$\widetilde r_i,\widetilde e_i$。以位置1为参考，两类损失可写成
\begin{align}
 \mathcal L_r&=-\sum_{i:C_i=1}
       \frac{\widetilde e_1}{\widetilde e_i}\log\widetilde r_i,\\
 \mathcal L_e&=-\sum_{i:C_i=1}
       \frac{\widetilde r_1}{\widetilde r_i}\log\widetilde e_i.
 \label{eq:rs-dla}
\end{align}
更新一个模型时，另一个模型提供校正权重；训练将两种估计互相联系，而线上只保留排序分数。

假设第三位估计检查分数仅为首位的五分之一，第三位点击在排序损失中就取得5倍权重。另一方面，如果第三位对象本来就被估计为很不相关，位置模型也会据此调整“低点击究竟来自哪里”的判断。列表softmax是训练参数化，不能不加条件地解释成每个对象的独立相关概率。

理论依赖检查、相关性与点击的分解及相应优化条件。两个网络能共同拟合日志，不保证任意数据下的可识别性；若相关性模型缺乏表达能力，或位置还通过信任直接改变点击，互估也可能偏离目标。与Regression-based EM的离线交替估计不同，DLA在同一训练循环中联合更新两个网络。

% Outline: 4.3.3.4.4
\paragraph{PAL：基于检查与点击分解的响应学习}
\label{sec:rs-pal}

普通CTR模型若把位置作为输入，服务排序时却不知道候选最终位置，只能填默认值，训练时学到的位置效应无法干净剥离。\term{PAL}把位置相关检查概率与查看后点击概率分开：ProbSeen只读取位置，pCTR读取用户、对象及其他特征，以两者乘积拟合日志点击\scite{rs-r92}

设$q_\eta(k)$为检查分支，$r_\theta(\bm x,i)$为响应分支，训练目标为
\[
 \ell_{\mathrm{BCE}}\!\left(C,q_\eta(k)r_\theta(\bm x,i)\right).
\]
检查分支与响应分支共同解释日志中的位置CTR差异。教学例中，首位与末位检查概率为0.8、0.2，同一商品查看后点击率为0.25，乘积分别为0.20、0.05。线上只使用$r_\theta$比较候选，无需提前猜一个位置，这也是分支分解相对“位置当输入特征”的服务优势。

结构成立依赖检查只由位置决定、查看后点击不再随位置变化等假设。若把$q$乘一个常数、$r$除以该常数仍保持合法范围，可能得到相同点击乘积，说明拟合本身不足以固定绝对尺度。原任务是推荐CTR，迁移到搜索时还要重新说明“查看后响应”与相关性标签之间的关系；PAL只剥离查看与响应，不处理信任点击带来的偏移，后者由仿射校正承接。

% Outline: 4.3.3.4.5
\paragraph{Affine Correction：基于仿射点击模型的排序学习}
\label{sec:rs-affine-correction}

PBM把点击写成查看乘相关性，隐含“不相关对象被查看后不会点击”；但用户可能因信任靠前结果而点击不相关对象，单纯除以检查概率（纯逆倾向只有缩放）不能去掉这种额外点击。\term{仿射校正}（affine correction）把位置条件下的点击概率写成相关性的缩放加偏移\scite{rs-s34}
\begin{equation}
 \mathbb P(C=1\mid i,k)=\alpha_k r_i+\beta_k,\qquad
 \widehat r_i=\frac{C-\beta_k}{\alpha_k}.
 \label{eq:rs-affine}
\end{equation}
请求条件在式中省略。若查看概率为$e_k$，查看后对相关、无关对象的点击概率分别为$\epsilon_k^+,\epsilon_k^-$，则
$\alpha_k=e_k(\epsilon_k^+-\epsilon_k^-)$，
$\beta_k=e_k\epsilon_k^-$。
需要$\alpha_k>0$及合法概率范围。

假设$\alpha_k=0.5,\beta_k=0.2$，点击时校正贡献为1.6，未点击时为$-0.4$。它们不是概率，但条件期望等于$r_i$，可进入相关性加权的排序目标。纯逆倾向只有缩放，无法消除$\beta_k$这项偏移。

在线使用学到的相关性分数，校正系数只需在训练和评价中按协议使用。系数估计错误、没有跨位置覆盖，或真实信任机制偏离仿射形式，都可能破坏校正。允许负训练贡献也要求使用与目标推导相容的损失，不能把贡献截到$[0,1]$后仍声称保持原无偏性；仿射模型仍是位置条件下的线性假设，更复杂的信任路径由后续结构或实验证据处理。

% Outline: 4.3.3.4.6
\paragraph{DR-LTR：基于偏好预测与加权残差的排序学习}
\label{sec:rs-dr-ltr}

低位逆倾向权重$1/e_k$可能很大，放大单条点击的方差；而一般双重稳健估计又要求实际处理可见，点击日志中用户是否查看通常不可见，无法按实际查看做补正。Oosterhuis的\term{DR-LTR}用各位置的预期作用代替不可见的实际查看，将偏好预测（回归基准）与点击残差结合\scite{rs-s35}

用式~\eqref{eq:rs-affine}的系数，令日志策略将对象$i$放到位置$k$的边际概率为$\mu(k\mid i)$，定义
$\rho_i=\sum_k\mu(k\mid i)\alpha_k$。
在固定请求的教学记法下，偏好估计具有如下形式：
\begin{equation}
 \widehat r_i^{\mathrm{DR}}
 =\widehat r_i^{\mathrm{reg}}
  +\frac1N\sum_{n=1}^{N}
   \frac{C_{n,i}-\alpha_{k_{n,i}}\widehat r_i^{\mathrm{reg}}
                     -\beta_{k_{n,i}}}
        {\widehat\rho_i}.
 \label{eq:rs-dr-ltr}
\end{equation}
未进入该次展示的对象约定$C=\alpha=\beta=0$。实现可以截断分母，并将对象估计与目标位置权重结合，形成可优化的排序准则。原文还给出与点击模型相容的偏好回归训练。

假设$\widehat r_i^{\mathrm{reg}}=0.4$，当前位置$\alpha=0.5,\beta=0.1$，估计$\widehat\rho_i=0.25$。一次点击的预测点击率为0.3，残差贡献为$(1-0.3)/0.25=2.8$；若只有这条教学记录，基准加残差为3.2。大于1并非计算错误，因为它是校正估计，不是单样本概率。多次记录平均后才用于估计目标。

双重稳健的范围需要完整陈述：位置和信任参数须正确；在此前提下，对每个对象，日志位置分布估计准确且截断不生效，或该对象的偏好回归准确，可以取得无偏估计。回归准确时，残差期望为零；位置分布准确时，残差补回回归误差。它不意味着任意两个网络只需一个正确，更不等同第~\ref{chap:rs-evaluation}~章的一般策略DR。与纯IPS相比，DR-LTR用回归基准压低残差项的方差，代价是多一套位置分布与偏好回归估计。

表~\ref{tab:rs-position-methods}汇总各方法在训练中改变的位置。它们可能在相容假设下组合，年份并不规定必经的替代顺序。
\begin{table*}[htbp]
 \centering\small
 \begin{tabularx}{\linewidth}{@{}lXXXX@{}}
 \toprule
 方法 & 倾向或参数来源 & 生成假设 & 训练校正与线上输出 & 主要失败条件\\
 \midrule
 Propensity SVM-Rank & 换位实验或可比较位置记录 & 查看乘相关性 & 点击名次损失加权；特征分数 & 小倾向、缺覆盖、信任作用\\
 Regression EM & 特征回归与位置参数交替估计 & 查看、相关性潜变量 & 估计倾向供排序训练；最终排序器 & 分解不唯一、回归失配\\
 DLA & 位置与排序模型互估 & 检查与相关性可分解 & 两类加权列表损失；排序分数 & 互估错误、优化与识别条件不足\\
 PAL & 位置分支与响应分支联合拟合 & 位置仅影响查看 & 乘积点击损失；查看后响应 & 绝对尺度、结构失配\\
 仿射校正 & 位置缩放及信任偏移估计 & 仿射点击模型 & 减偏移再缩放；相关性分数 & 系数错误、零缩放或缺覆盖\\
 DR-LTR & 点击模型、日志位置分布及偏好回归 & 正确位置与信任参数 & 基准加残差；训练后排序器 & 超出双重稳健条件、截断偏差\\
 \bottomrule
 \end{tabularx}
 \caption{从带位置点击学习候选评分的方法。相同对象在不同位置的可比记录有助于区分查看和相关性；没有覆盖或额外约束时，结构分支本身不能证明分解正确。}
 \label{tab:rs-position-methods}
\end{table*}

% Outline: 4.3.3.5
\subsubsection{流行度偏差下的评分学习}
\label{sec:rs-popularity-bias}

热门物品可能因为更多曝光而积累点击，也可能因为销量展示影响选择；优质内容及个性化匹配同样可能使其热门。因而“减去热门程度”不能自动恢复偏好。下面两种方法分别改变反事实评分与表示分解，均需保留其结构假设，并与真实干预证据区分。

% Outline: 4.3.3.5.1
\paragraph{MACR：基于反事实评分的流行度纠偏}
\label{sec:rs-macr}

热门物品因更多曝光而积累点击，普通匹配模型会把这部分流行度优势学进用户—物品分数，事后再手工“减去热门程度”缺乏统一依据。\term{MACR}在用户—物品匹配分支之外，增加用户和物品各自的边际分支，以同一交互标签联合训练主评分和辅助任务\scite{rs-r93}

\textbf{模型结构。}第一步，匹配分支产出用户—物品匹配分数$m_{u,i}$；用户、物品边际分支分别sigmoid输出$a_u,b_i$，事实组合为$m_{u,i}a_ub_i$。第二步，服务时把匹配分数替换为参考常数$c$，构造反事实组合，再取差得到排序分数：
\begin{equation}
 s_{\mathrm{MACR}}(u,i)=(m_{u,i}-c)a_ub_i.
 \label{eq:rs-macr}
\end{equation}
这里保留了边际分支；差值是候选排序分数，不是校准点击概率。

教学例取$a_u=0.8$，热门候选A的$m=0.6,b=0.9$，候选B的$m=0.8,b=0.4$。事实组合分别为0.432、0.256，A领先；若参考$c=0.5$，校正分数变为0.072、0.096，B领先。参考状态使匹配较弱但边际较强的候选受到更大扣除，同时也说明$c$会影响排序，不能隐藏为无关常数。

论文以因果图解释匹配路径和直接路径，模型可接MF或LightGCN等骨干。但边际分支由交互标签学到什么，仍受曝光、质量和其他未建模因素影响。分支命名及反事实计算并不单独证明真实流行度原因已被识别，效果应按人群、流行度区间及对照条件检查；与DICE的表示分离相比，MACR在评分端做反事实扣除。

% Outline: 4.3.3.5.2
\paragraph{DICE：基于兴趣与从众分离的流行度纠偏}
\label{sec:rs-dice}

MACR在评分端扣除流行度，却仍把兴趣与从众混在同一匹配表示里；要从源头分开，需要让两套表示各自解释不同因素。\termalias{recommendation-dice}{DICE}为用户和物品分别学习兴趣、从众两套表示，用流行度关系构造不同辅助监督\scite{rs-r94}

\textbf{模型结构。}第一步，分别学习兴趣表示$\bm h_u^{\mathrm{int}},\bm v_i^{\mathrm{int}}$与从众表示$\bm h_u^{\mathrm{con}},\bm v_i^{\mathrm{con}}$。第二步，最终评分保留两项相加：
\begin{equation}
 s(u,i)=
 \langle\bm h_u^{\mathrm{int}},\bm v_i^{\mathrm{int}}\rangle+
 \langle\bm h_u^{\mathrm{con}},\bm v_i^{\mathrm{con}}\rangle.
 \label{eq:rs-dice}
\end{equation}
它试图让两个分支解释不同因素，而不是服务时默认删除从众项。

第三步按流行度关系给不同成对信号：若正例$i$比采样负例$j$更冷门，用户仍选择$i$，在模型假设下，兴趣分支应更偏向$i$，从众分支却应更偏向热门$j$；总点击分数仍应满足$i$领先。若正例更热门，只能为总分和从众提供相应约束，不应强迫兴趣也更偏向正例。两类样本因此产生不同的成对训练信号，再与表示分离损失联合学习。

例如冷门水壶的兴趣、从众得分为1.2、0.1，热门水壶为0.7、0.4，总分分别为1.3、1.1；这个教学分解允许兴趣抵消热门优势。其解释依赖两因素结构、独立性、加性评分及流行度代理。未点击仍可能因为未曝光，数据中的分离效果也不足以证明潜在原因唯一。

% Outline: 4.3.3.6
\subsubsection{隐私约束下的评分学习}
\label{sec:rs-private-scoring}

用户记录不能集中汇集时，评分任务没有消失，但可交换的信息变了。联邦学习规定客户端如何协作更新模型；差分隐私限制输出对单条或单个用户数据的敏感程度。二者职责不同，不能以“只上传梯度”代替隐私分析。这里关注本地图、更新和服务表示分别留在哪里。

% Outline: 4.3.3.6.1
\paragraph{FedPerGNN：分散交互图上的协同评分}
\label{sec:rs-fedpergnn}

协同过滤依赖用户—物品交互图，集中汇集图意味着把个人历史交给服务器，与隐私约束冲突；联邦学习要求原始记录留在设备，只交换受保护的更新。\term{FedPerGNN}从用户设备上的交互子图学习用户、物品表示，以已观测评分监督预测；学习服务器聚合受保护更新，可信第三方帮助发现匿名邻居\scite{rs-r122}

\textbf{模型结构。}第一步本地训练：设备在自己的交互子图上完成GNN传播与评分预测。第二步保护更新：客户端为伪交互物品加入更新，减弱稀疏更新索引暴露真实交互的风险，再裁剪梯度并添加拉普拉斯噪声，才上传给服务器聚合。第三步图扩展：可信第三方匹配加密物品标识，分发共同交互者的匿名表示，为本地提供高阶关系。

例如设备A和B都评价过电影$i$，原始评分及完整历史留在各自设备。第三方看到可匹配的加密标识，为两侧分发匿名邻居信息；学习服务器收到经保护的模型更新，再聚合为新的共享参数。本地用户表示由本地记录及邻居关系更新。原方案训练后还将用户表示上传到学习服务器以提供个性化服务，因此不能将整条流程表述为“任何用户表示都不离开设备”。

隐私分析需分别覆盖真实交互索引、数值梯度、匿名邻居和最终表示。服务器与第三方的信任分工是条件；多轮更新还会累计隐私消耗。更强噪声、更小裁剪阈值或更多伪物品会影响准确度、通信与计算，不能仅凭原始数据保存在本地就断言没有泄漏。

% Outline: 4.3.3.7
\subsubsection{噪声与恶意反馈下的评分学习}
\label{sec:rs-robust-scoring}

误点、错误标签、虚假交互和恶意客户端更新进入模型的路径不同。参数扰动训练可以增强某种局部稳定性，模型投毒实验则检查攻击者能否改变共享参数。两类工作分别回答防护目标和风险入口，也都不能替代曝光、位置或流行度纠偏。

% Outline: 4.3.3.7.1
\paragraph{APR：参数扰动下的稳健排序训练}
\label{sec:rs-apr}

BPR只在当前参数上优化成对分差，正负候选分差很小时，小幅参数变化（数值噪声、在线漂移、部分更新）就可能翻转次序，而训练不会主动提防这种局部脆弱方向。\term{对抗个性化排序}（adversarial personalized ranking, APR）在BPR目标上寻找范数受限的增损参数扰动，再同时降低原始与扰动后的损失\scite{rs-r123}

\textbf{模型结构。}第一步找最坏扰动：令$\mathcal L_{\mathrm{pair}}$只表示式~\eqref{eq:rs-bpr}的成对数据损失，在半径$\epsilon$的范数球内最大化扰动损失，
\begin{align}
 \bm\delta^\star&=\argmax_{\lVert\bm\delta\rVert_2\leq\epsilon}
       \mathcal L_{\mathrm{pair}}(\bm\theta+\bm\delta),\\
 \mathcal L_{\mathrm{APR}}
   &=\mathcal L_{\mathrm{pair}}(\bm\theta)
     +\lambda_{\mathrm{adv}}
       \mathcal L_{\mathrm{pair}}(\bm\theta+\bm\delta^\star)
       +\frac{\lambda}{2}\lVert\bm\theta\rVert_2^2.
 \label{eq:rs-apr}
\end{align}
第二步联合更新：原始损失与最坏扰动损失一同回传，使参数在邻域内都保持次序。原文以MF实现对抗矩阵分解，先用BPR初始化，再以梯度方向近似最坏扰动，并交替执行扰动构造和参数更新。

在一维教学例中，用户因子为1，正负物品因子为0.51、0.50，原分差仅0.01。把正物品减0.01、负物品加0.01，总扰动范数约0.0141，就会使分差变成$-0.01$。APR让训练关注这种局部脆弱方向，稳健更新后仍按原评分器比较候选，不在服务时施加恶意扰动。

评价应同时报告普通排序质量、规定扰动下的质量和额外训练成本。参数范数球是明确选择的威胁模型，不等同用户刷量、伪造日志或任意客户端投毒；不能把局部稳健性扩大成对未知攻击的保证。它与FedRecAttack的区别在于：APR扰动训练中规定的参数邻域以增强稳健性，FedRecAttack研究恶意上传更新以破坏共享模型。

% Outline: 4.3.3.7.2
\paragraph{FedRecAttack：恶意更新下的评分风险诊断}
\label{sec:rs-fedrecattack}

联邦训练把部分更新交给服务器聚合，默认客户端上传的是改善各自误差的正常更新，也就需要检查上传更新是否可信。\term{FedRecAttack}研究受控客户端怎样利用公开交互近似用户表示，再上传目标导向的恶意梯度，推高指定物品的共享表示和排名\scite{rs-r124}它是模型投毒研究，不是一个服务端防御算法。

\textbf{模型结构。}第一步，受控客户端利用公开可见的交互近似目标用户表示。第二步，构造目标导向的恶意更新：正常客户端的更新改善各自评分误差，受控客户端的更新却集中增加目标商品的得分。第三步，服务器按比例聚合后，恶意方向进入共享参数，原本没有对目标商品表现兴趣的用户也可能看到它上升；随后产生的曝光又可能进入新日志，形成自我强化。教学分析应沿“公开交互—近似用户表示—恶意更新—共享物品参数—候选次序”检查影响，并明确受控客户端比例、可见信息和更新限制。

只报告目标物品命中率不够，还要看普通用户的排序损失、其他物品受损情况及整体质量，以免把破坏所有排序误当成隐蔽攻击成功。APR扰动的是训练中规定的参数邻域以增强稳健性，FedRecAttack改变的是上传更新的完整性；FedPerGNN保护记录或更新隐私，也不自动保证参与者诚实。

% Outline: 4.3.4
\subsection{评分应用与维护}
\label{sec:rs-scoring-maintenance}

模型输出后，还要确定它在什么决策中使用、在哪个阶段执行，以及多久以后仍然有效。概率校准、阶段衔接、时间更新和场景适配可以共同作用于同一评分器。数据撤回则改变允许保留的学习依据，需要清除已有影响，不能只靠加入更新的行为解决。

% Outline: 4.3.4.1
\subsubsection{分数解释、校准与不确定性}
\label{sec:rs-score-calibration}

概率校准要求预测值与同类机会中的响应频率对应。例如所有预测0.1的曝光，长期点击频率应接近0.1。总体平均正确仍可能掩盖页面或人群内的偏差；按分数分箱及关键场景检查，才有助于发现混排和价值计算中的系统误差。单调校准可以保持多数次序，却不能把不同标签定义的分数变成统一效用。

不确定性回答的是“这个预测获得多少证据支持”。两个候选都预测0.1，一个附近有大量成熟样本，另一个来自新内容区域，其可靠程度可能不同。置信区间、模型间分歧、表示密度等各有假设；概率值靠近0或1并不自动意味着模型有充分证据。发生后价值、成本及决策规则还需与概率共同使用。

% Outline: 4.3.4.1.1
\paragraph{PRU：预测相关样本的不确定性估计}
\label{sec:rs-pru}

在点击稀少、正负类别重叠的数据中，靠近任意训练样本未必表示预测可靠；而深度集成、蒙特卡洛采样等不确定性方法要为每个候选重复多次前向，代价过高。\term{预测相关不确定性}（predictive relevance uncertainty, PRU）先按类别对训练损失拟合高斯混合，选取低损失的预测相关样本，再用这些样本的隐藏表示拟合各类高斯分布\scite{rs-r109}

\textbf{模型结构。}第一步筛样本：按类别对训练损失拟合高斯混合，保留低损失的预测相关样本，而非全部训练点。第二步建密度：用所选样本的隐藏表示拟合各类高斯分布，模型训练结合谱归一化与双侧Jacobian正则，鼓励隐藏距离保留输入变化，减少表示塌缩对密度判断的干扰。第三步算不确定性：设前向计算得到$\bm h$，所选可靠样本形成类别均值$\bm\mu_y$、协方差$\bm\Sigma_y$及混合权重$\omega_y$，不确定性依据可由
\[
 q(\bm h)=\sum_y\omega_y
   \mathcal N(\bm h;\bm\mu_y,\bm\Sigma_y)
\]
的低密度给出，例如使用负对数密度作为排序量。线上只需一次模型前向及密度计算，无须对每个候选重复采样许多网络。

两个pCTR同为0.1的候选中，A落在预测相关样本的高密度区域，B远离这些区域，可以对B采用更保守回退或在允许范围内安排探索。阈值仍要用验证数据确定。PRU输出不是点击概率，也不是具有覆盖保证的置信区间；高斯假设、类别样本不足和表征失真都会影响它。探索见第~\ref{chap:rs-presentation}~章，集合风险控制见第~\ref{chap:rs-constraints}~章。

% Outline: 4.3.4.2
\subsubsection{召回与排序的阶段衔接}
\label{sec:rs-stage-connection}

召回要从全目录取回对象，粗排要以较小成本筛选较多候选，精排则能给少量候选更多计算。重排还会处理列表关系。假设教学流程把十万候选缩至两千、再缩至一百，前阶段误删的对象通常无法由后阶段恢复；只在已经进入精排的样本上测精排质量，会遗漏这部分损失。

阶段可以共享表示、共同训练或使用教师蒸馏。例如教师对同一候选产生目标分布$q_T$，学生最小化与$q_T$的交叉熵或KL散度，使粗排保留精排认为好的对象。但教师只见过筛选后的候选，蒸馏也会继承这一覆盖限制。语言模型重排放在后阶段时，还需把调用次数、窗口处理和教师输出错误计入总成本。

下面四个方案分别调整用户编码共享、物品复用、跨阶段监督和生成—评分计算。它们仍保留候选与分数接口，不等同直接产生多物品展示列表。

% Outline: 4.3.4.2.1
\paragraph{OneTrans-V2：跨阶段共享与联合学习}
\label{sec:rs-onetrans-v2}

召回、粗排、精排分别编码相同历史，会重复消耗计算，三套独立用户状态还可能彼此不一致，粗排因此滤掉精排认为好的对象。\term{OneTrans-V2}让召回、粗排和精排读取共享用户骨干，同时保留各自特征单元及输出\scite{rs-r43}

\textbf{模型结构。}第一步共享骨干：三个阶段共用同一用户表示计算，召回采用第~\ref{sec:rs-dcgr}~节的DCGR生成物品标识；粗排、精排计算给定候选的响应分数。第二步联合训练：三个任务共同训练用户表示，精排再向粗排蒸馏，教师logit在蒸馏项中停止梯度。第三步隔离泄漏：阶段间和不同曝光的特征单元隔离，避免某次响应或后阶段候选向其他样本泄漏。

停止教师梯度只约束蒸馏路径，不表示精排任务不能更新共享骨干。共同任务损失仍由各自样本回传。原OneTrans的字段—历史计算见第~\ref{sec:rs-onetrans}~节，区别在于OneTrans统一评分内部计算，V2把共享扩到召回—粗排—精排三个阶段。

教学窗口内有请求$t_1,t_2$，中间新增行为$b$。$t_1$的三个阶段只读取前缀$H_{u,t_1}$，$t_2$读取$H_{u,t_1}$加上当时已到达的$b$；不能因两次请求合并训练，就让$t_1$看到$b$。追加式历史和因果掩码允许摊薄重复编码，缓存还要随模型、事件及特征版本维护。共享没有省去阶段执行和候选截断。

评价应同时记录各阶段质量、对精排高价值对象的保留率及总请求成本，区分稀疏专家扩容、参数化与服务优化各自作用。截至2026年10月5日，本工作按2026年9月预印本使用，其实验范围不能被扩大为所有业务的普遍结论。

% Outline: 4.3.4.2.2
\paragraph{IntTower：可复用物品计算的交互粗排}
\label{sec:rs-inttower}

粗排既想保留双塔物品编码可预存、离线更新的优势，又希望获得比一次单向量点积更细的交互；点积把所有字段压进一个方向，匹配粒度粗。\term{IntTower}在两塔内使用轻量字段注意，将用户各层表示和物品末层表示投影为多头，再作最大相似度聚合\scite{rs-r60}

\textbf{模型结构。}第一步塔内字段注意：用户塔各层、物品塔末层先做轻量字段注意，再投影为多头表示。第二步头间最大匹配：令用户第$\ell$层第$h$头为$\bm u_{\ell,h}$，物品头为$\bm v_g$，交互分数为
\begin{equation}
 s(u,i)=\sum_{\ell}\sum_{h=1}^{H}
           \max_{1\leq g\leq H}
               \langle\bm u_{\ell,h},\bm v_g\rangle.
 \label{eq:rs-inttower}
\end{equation}
物品各头可以合并存为一个结构化向量，服务时仍需恢复头的分组和最大匹配，而不是退化成一次普通点积。

在单层两头教学例中，用户的两个头分别为$(1,0)$与$(0,1)$。候选A的两个头为$(0.8,0.2)$与$(0.1,0.6)$，两次最大值为0.8、0.6，总分1.4。候选B的两个头为$(0.5,0.5)$与$(0.4,0.4)$，总分1.0。计算需要各头间比较，而不是只算拼接向量的一次普通点积；若还有用户层，继续累加对应分数。

将聚合分数经sigmoid变成响应预测后，训练结合点击交叉熵、用户—正物品表示的对比正则及参数正则。物品侧预存、在线用户计算和逐候选交互共同决定时延。模型的原任务是粗排，保留双塔结构不意味着它自动成为普通单向量近邻索引，也不意味着召回、精排已联合训练。

% Outline: 4.3.4.2.3
\paragraph{HCCP：跨阶段一致性与候选覆盖}
\label{sec:rs-hccp}

只模仿精排已曝光样本，粗排就很难学习自己早先滤掉的对象：被截断的物品没有点击标签，模型也从不比较它们，覆盖盲区由此形成。\term{HCCP}从精排候选、粗排过滤部分、批内物品和目录采样构造不同层次的候选，再组合跨阶段顺序一致性与对比训练\scite{rs-r61}

\textbf{模型结构。}第一步构造多层次候选池：精排候选、粗排过滤部分、批内物品、目录采样分别进入训练集。第二步跨阶段一致性：精排序列结合曝光和点击反馈重新组织，全局列表损失对齐较大候选空间，局部列表损失强调曝光空间内的次序。第三步难度区分：带间隔的对比目标区分不同来源的候选难度。

例如$a,b$进入精排且$a$点击，$c$被粗排截断，$d$来自同批其他用户，$e$来自目录随机采样。$a,b$可以提供曝光响应及局部顺序，$c$有上游筛选信息却没有当前真实点击标签，$d,e$主要提供额外竞争。训练时不能把这五者都解释为同样可靠的已观察正负反馈。

线上仍由粗排模型输出分数并截断。应检查其对精排偏好对象的保留率、被过滤及长尾候选的覆盖，以及后续实际响应。不同负例难度是训练先验，扩大候选空间并不自动证明已去除选择偏差。HCCP改变监督组织，OneTrans-V2改变共享骨干，二者是不同坐标。

% Outline: 4.3.4.2.4
\paragraph{UniR\textsuperscript{2}：生成轨迹与多目标评分共享}
\label{sec:rs-unir2}

生成物品标识（召回）与评估给定候选（排序）需要的信息并不相同，两套独立模型又会重复编码用户、状态不一致。UniR$^2$把用户条件、标识生成轨迹和候选字段放进一条序列，使用两类查询及不同可见范围\scite{rs-r74}

\textbf{模型结构。}第一步双查询与掩码：生成查询读取用户前缀和过去代码，按自回归逐个产出标识；评分查询读取画像、完整生成轨迹及候选字段，不再直接读取长历史。候选字段内部可以相互读取，生成路径却不能看到尚未生成的代码。第二步双损失：生成路径使用代码负对数似然，评分路径使用各响应标签的联合损失；目标仍是候选及其响应分数，没有直接确定最终展示列表。

例如召回依次生成代码$(c_1,c_2,c_3)$，取得候选$i$；评分阶段再补入价格、作者及交叉统计等字段，结合这一生成轨迹形成多目标预测。

第三步参数隔离：两条路径共享基础注意权重，使用不同前馈网络；评分侧通过LoRA和停止梯度隔离对基础生成计算的更新。可概括为评分投影
$\operatorname{sg}(\bm W\bm x)+\bm B\bm A\bm x$，
其中$\operatorname{sg}$阻断相应基础路径梯度，低秩支路保留评分适应能力。缓存复用的是相同用户条件与生成轨迹，不能据“单次前向”就声称全部参数共同优化或所有信息双向流动。

% Outline: 4.3.4.3
\subsubsection{分布变化下的评分更新}
\label{sec:rs-score-drift}

一次促销会改变价格和候选，一次页面改版会改变查看概率，用户兴趣也可能随时间变化。误差上升后，应该先区分输入更新不及时、标签尚未成熟、条件响应变化及业务目标改变。前两章分别处理需求依据和目录维护，本节关注它们进入评分后是否仍满足训练条件。

如果所有候选概率近似按同一比例偏移，重新校准可能足以修复价值计算；如果用户—物品次序改变，则只调整校准曲线不够。更新可采用滚动窗口、增量学习或重新训练，但需防止短窗口丢掉稳定的低频偏好，也要检查旧阶段输出与新阶段输入是否匹配。

评价不能只看更新后的单日指标。应记录变化发生前后误差、恢复时间、历史能力保留和维护成本，并使用当时可得的信息重放。更快更新还可能更快吸收误点或污染；时间新鲜度本身不是收益保证。

% Outline: 4.3.4.4
\subsubsection{多场景条件下的评分适配}
\label{sec:rs-multiscenario}

同一用户与商品出现在首页、购物车或购后页时，候选上下文和点击基线都可能不同。场景条件描述“在哪里作预测”，响应任务描述“预测什么”；不能把场景数与响应头数混为一谈。相较第~\ref{sec:rs-emcdr}、\ref{sec:rs-conet}~节的跨域迁移，这里主要考虑共享目录下同时存在的页面分布差异。

% Outline: 4.3.4.4.1
\paragraph{STAR：场景共享与专属评分}
\label{sec:rs-star}

同一用户与商品出现在首页、购物车或购后页时，完全独立建模浪费共享信息，完全共享又可能掩盖场景差异（候选上下文与点击基线不同）。\term{STAR}由场景标识选择专属参数，与共享中心参数组合；分区归一化则为不同场景保留相应统计和变换\scite{rs-r57}

\textbf{模型结构。}第一步中心—专属组合：一层有效权重为$\bm W_{\mathrm c}\odot\bm W_d$（共享中心权重逐元乘场景专属权重），偏置由中心与场景偏置相加，随后作激活；模型还可以通过辅助分支融合场景信息。第二步分区归一化：为不同场景保留相应统计和变换，避免大场景统计主导小场景。

用单元教学例看梯度分工：共享权重$w_{\mathrm c}=2$，首页专属权重0.5，购物车专属权重1.5，输入均为0.4且偏置为零，则两个场景logit为0.4、1.2，概率约为0.599、0.769。首页样本更新中心和首页参数，购物车样本更新中心和购物车参数；两个场景通过中心权重交流，专属分支则保留区别。实际网络还会执行场景内归一化及多层变换。

应同时比较整体与场景内误差、校准和资源，防止大场景平均值掩盖小场景退化。该方法需要已知场景和足够样本，不是未知场景或持续漂移的通用解；预测任务之间的共享仍由MMoE等结构另行处理。

% Outline: 4.3.4.4.2
\paragraph{PEPNet：场景与个体先验的条件评分}
\label{sec:rs-pepnet}

场景内部，不同用户、对象和作者也可能需要不同的计算贡献；为每个用户独立保存整套网络在参数与存储上都不可行。\term{PEPNet}通过嵌入个性化网络EPNet和参数个性化网络PPNet生成门控，以轻量缩放实现个性化，而不是为每个用户独立保存整套网络\scite{rs-r58}门控基本单元采用两层网络，末端以$2\sigma(\cdot)$产生围绕1的缩放。

\textbf{模型结构。}第一步嵌入个性化（EPNet）：将场景先验与停止梯度的通用嵌入拼接，生成逐维门控，再乘到共享嵌入上。第二步参数个性化（PPNet）：将用户、对象、作者先验与停止梯度的EPNet输出结合，为各任务、各层生成隐藏单元门控。主预测路径仍能更新共享表示，停止梯度阻止的是门控输入这条旁路，并非冻结所有嵌入。

若点击任务某层表示为$(1,0.5)$、门控为$(1.4,0.6)$，缩放后为$(1.4,0.3)$；观看任务门控取$(0.8,1.2)$，相应结果为$(0.8,0.6)$。后续任务头分别拟合自己的响应标签，联合损失更新门控和网络。与MMoE选择专家不同，这里直接调整表示和隐藏单元的贡献；得到多个预测后，最终目标权衡仍由决策规则规定。

% Outline: 4.3.4.5
\subsubsection{数据撤回条件下的评分维护}
\label{sec:rs-unlearning-maintenance}

用户撤回历史后，删除原始日志并不能清除已经进入嵌入、图关系和模型参数的影响。撤回任务要明确删除哪些记录及其派生结果，再取得只依赖允许保留数据的评分器。它与需求变化时吸收新数据不同，也不由差分隐私自动完成。

% Outline: 4.3.4.5.1
\paragraph{RecEraser：分片重训与自适应评分聚合}
\label{sec:rs-receraser}

用户撤回历史后，删除原始日志并不能清除已经进入嵌入、图关系和模型参数的影响；全量重训是直接参照，但每次撤回都执行代价很高。\term{RecEraser}先按用户、物品或交互相似性构造均衡分片，在各片独立训练子模型，再学习注意聚合用户和物品表示\scite{rs-r125}

\textbf{模型结构。}第一步分片：按用户、物品或交互相似性构造均衡分片，各片独立训练子模型。第二步自适应聚合：用户表示为$\bm h_u=\sum_j a_{u,j}\widetilde{\bm h}_{u,j}$，物品表示也由相应子模型表示聚合；必要投影将子模型空间对齐，再按选定评分函数输出。

撤回时若用户删除一条电影交互，该记录只属于某个固定分片，就从该片删除并重训受影响子模型，同时重训聚合部分。其他独立子模型不必因这一条记录全部重训。分片过多会破坏协同关系，过少则让一次删除涉及较多训练，因此分片数需要按保留质量与删除成本选择。

维护范围不能止于子模型文件。分片若依赖包含该记录的预训练表示，完整流水线与“从未使用该记录重新开始”的参照仍可能不同；邻域、缓存、聚合器及部署后的派生表示也要检查。固定分片框架下的删除性质与完整训练流程的等价性应分别评价，后续第~\ref{chap:rs-evaluation}~章进一步讨论用户删除及协同影响。

% Outline: 4.3.5
\subsection{评分效果评价}
\label{sec:rs-scoring-evaluation}

响应概率、数值价值和排序次序各需要相应检查。点击概率可用交叉熵、Brier分数及分组校准，数值价值可用MAE、RMSE及关键区间误差；AUC衡量正负区分，NDCG等截断指标检查头部次序。AUC不检验概率尺度，低平均回归误差也不保证前几名正确。增量排序则需处理与对照证据，不能直接以普通点击NDCG代替。

这些指标都受标签窗口、采样和观测机制影响。比较评分模型时，固定请求时信息、候选和预算；比较阶段方案时，记录候选截断、教师覆盖及全流程成本。要区分新增历史、更多字段、更多计算和更及时更新各自的贡献，消融比较才有可解释的对象。

即使以上检查通过，评分器仍只是展示决策的一部分。最终列表的重复、位置、候选之间的竞争与多方约束，可能使单对象预测改善没有转化为实际收益。第~\ref{chap:rs-evaluation}~章建立完整离线、策略和实验协议，下一章则研究怎样把这些分数组织成用户实际看到的展示。

% Outline: 4.4
\section{本章小结}
\label{sec:rs-ranking-summary}

给定候选后的比较必须从目标量开始。响应预估描述在明确条件下发生什么，价值预估还要结合结果的价值；偏好排序学习相对次序，增量排序则比较推荐与不推荐造成的结果差异。相同评分架构可以服务这些任务，但标签、损失和识别证据不能随名称一起互换。

字段交互、候选相关历史与语义内容提供评分依据，分离模块或联合骨干决定这些依据何时相遇。评分内部联合计算和跨阶段共享具有不同范围，复用又取决于时间与信息可见性。可靠学习还需要理解观测选择、位置及流行度作用，明确隐私保护与恶意更新的职责，并在数据撤回时处理已经学入的影响。

概率校准、证据可靠程度、排序质量和最终收益需要分别检验。即使每个候选都预测准确，也不能据此确定它们应该怎样并列、排位或反复出现。下一章从单对象分数转向展示列表与交互，将“谁更合适”进一步落实为“用户这一次会看到什么”。
